% concatenated from main.tex
%\documentclass[draftclsnofoot, onecolumn, comsoc, 12pt]{IEEEtran} 
\documentclass[10pt, twocolumn, comsoc]{IEEEtran}
% \documentclass[conference]{IEEEtran}

\usepackage[T1]{fontenc}

\newcommand{\latexiie}{\LaTeX2{\Large$_\varepsilon$}}
\def\argmax{\mathop{\rm arg\,max}}
\def\argmin{\mathop{\rm arg\,min}}
\usepackage{graphicx,epsfig}
\usepackage[noadjust]{cite}
\usepackage{mcite}
\usepackage{amsfonts,helvet}
\usepackage{fancyhdr}
\usepackage{threeparttable}
\usepackage{epsf,epsfig}
\usepackage{amsthm}
\usepackage{amsmath}
\usepackage{siunitx}
\usepackage{amssymb}
\usepackage{stfloats}

\usepackage{dsfont}
% \usepackage{subfigure}
\usepackage[caption=false,font=footnotesize]{subfig}
\usepackage{color}
% \usepackage{algorithm}
% \usepackage{algpseudocode}
% \usepackage{algcompatible}
\usepackage{enumerate}
\usepackage{gensymb}
%\usepackage{ulem}
\usepackage{cancel}
%\usepackage{mathrsfs}
%\usepackage{pslatex}
%\usepackage{mathptmx}
%\usepackage{graphicx}
%\usepackage{times}
%\usepackage{txfonts}
%\usepackage{pxfonts}
%\usepackage{balance}
\usepackage{lipsum}
\usepackage{mathtools}
\usepackage{cuted}
\usepackage{bbm}
\usepackage[linesnumbered,ruled]{algorithm2e}
\usepackage[colorlinks=true, linkcolor=blue]{hyperref}

% \usepackage{algorithm}
\usepackage{algorithmic}

\newcommand{\eref}[1]{(\ref{#1})}
\newcommand{\fig}[1]{Fig.\ \ref{#1}}
\newtheorem{theorem}{Theorem}
\newtheorem{acknowledgement}{Acknowledgement}
\newtheorem{algo}{Algorithm}
\newtheorem{axiom}{Axiom}
\newtheorem{case}{Case}
\newtheorem{claim}{Claim}
\newtheorem{conclusion}{Conclusion}
\newtheorem{condition}{Condition}
\newtheorem{conjecture}{Conjecture}
\newtheorem{corollary}{Corollary}
\newtheorem{criterion}{Criterion}
\newtheorem{definition}{Definition}
\newtheorem{example}{Example}
\newtheorem{exercise}{Exercise}
\newtheorem{lemma}{Lemma}
\newtheorem{notation}{Notation}
\newtheorem{problem}{Problem}
\newtheorem{proposition}{Proposition}
\newtheorem{remark}{Remark}
\newtheorem{solution}{Solution}
\newtheorem{summary}{Summary}
\newtheorem{assumption}{Assumption}

\usepackage{eucal}
% \usepackage{geometry}
\usepackage{booktabs}

\usepackage{multirow}

\setcounter{page}{1}
%\newcounter{algo}
%\newcounter{proposition}
\setcounter{proposition}{0}

%\def\thefootnote{}

\newcommand{\C}{\mathbb{C}}%{{\mbox{\rm $\scriptscriptstyle ^\mid$\hspace{-0.40em}C}}} %o
\newcommand{\R}{\Re}%{{\mbox{\rm $\scriptscriptstyle ^\mid$\hspace{-0.40em}C}}} %o
\newcommand{\Z}{\mathbb{Z}}%{{\mbox{\rm $\scriptscriptstyle ^\mid$\hspace{-0.40em}C}}} %o 

%new commands for math font
\newcommand{\bb}[1]{\mathbb{#1}}
\newcommand{\bF}[1]{\mathbf{#1}}
\newcommand{\Cal}[1]{\mathcal{#1}}
\newcommand{\Frak}[1]{\mathfrak{#1}}

\renewcommand\qedsymbol{$\blacksquare$}

% \SetKw{KwBy}{by}

\begin{document}

\title{Fronthaul Compression for Uplink Cloud-RAN with Finite-Alphabet Inputs:\\ A Reverse Mercury/Waterfilling Approach}


    \author{Subin~Shin, Jaehoon~Lee, Seok-Hwan~Park, and Jeonghun~Park
\thanks{
This work was supported in part by the Institute of Information \& Communications Technology Planning \& Evaluation (IITP) grant funded by the Korea government (MSIT) (No. RS-2024-00395824, Development of Cloud virtualized RAN (vRAN) system supporting upper-midband; No. RS-2024-00404972, Development of 5G-A vRAN Research Platform; No. RS-2024-00435652, 6GARROW: 6G AI-native integrated RAN-Core networks), and in part by the 6GARROW project funded by the Smart Networks and Services Joint Undertaking (SNS JU) under the European Union’s Horizon Europe research and innovation programme (Grant Agreement No. 101192194).
S. Shin, J. Lee and J. Park are with the School of Electrical and Electronic Engineering, Yonsei University, Seoul 03722, South Korea (e-mail:{\texttt{ sbshin@yonsei.ac.kr, jeje5934@gmail.com, jhpark@yonsei.ac.kr}}). 
S.-H. Park is with the School of Electrical Engineering, Hanyang University, Ansan 15588, South Korea (e-mail:{\texttt{ seokhwanpark@hanyang.ac.kr}}).
}
}

 
\maketitle \setcounter{page}{1} 
\begin{abstract}
% The cloud radio access network (C-RAN) architecture mitigates inter-cell interference (ICI) in dense cellular networks by aggregating observations from distributed remote units (RUs) at a centralized unit (CU) for joint signal processing, but its capacity-limited fronthaul links require efficient signal compression.
% A widely adopted approach is transform-compress-forward, in which each RU first transforms its received signal into a lower-dimensional signal and then independently quantizes the resulting components.
% Its companion is bit allocation, which decides how the finite bit budget is distributed across the transformed components, and which has been comparatively
% under-explored.
% In this paper, we revisit the bit allocation problem under finite-alphabet inputs and propose a framework that explicitly accounts for their saturation behavior while taking the finite-alphabet generalized mutual information (GMI) induced by the linear minimum mean squared error (LMMSE) receiver.
% Applying the I-MMSE relation on the bit allocation side of the communication chain yields a closed-form fixed-point update whose converged form decomposes into a vessel height, a shared water level, and a finite-alphabet mercury level that explicitly accounts for the saturation behavior of finite-alphabet constellations.
% The resulting allocation, termed reverse mercury/waterfilling (RMWF), generalizes classical reverse waterfilling (RWF) through this saturation correction, with the classical form arising as a Gaussian high-SNR special case of our framework.
% Numerical results demonstrate that RMWF substantially outperforms both uniform and Gaussian-based allocations under tight fronthaul budgets and remains robust under antenna scaling where competing schemes degrade.
% The cloud radio access network (C-RAN) architecture mitigates inter-cell interference (ICI) by aggregating observations from distributed remote units (RUs) at a centralized unit (CU) for joint signal processing, but limited fronthaul capacity requires efficient signal compression. 
% A widely adopted approach is transform-compress-forward, in which each RU transforms its received signal and independently quantizes the resulting coefficients, with bit allocation determining how a finite bit budget is distributed across them.
% For Gaussian sources, the rate-distortion bit allocation is classically described by reverse waterfilling (RWF), but practical systems use finite-alphabet symbols whose mutual information saturates at $\log_2M$ rather than growing unboundedly with SNR.
% On the power allocation side, this asymmetry has been resolved by mercury/waterfilling (MWF), in which a saturation-aware mercury level augments classical waterfilling. 
% On the bit allocation side—precisely the operating regime of C-RAN fronthaul compression—an explicit finite-alphabet, MI-aware counterpart is still lacking.
% To fill this gap, we formulate bit allocation as the maximization of the post-linear minimum mean squared error (LMMSE) finite-alphabet generalized mutual information (GMI) at the CU. 
% Applying the I-MMSE relation on the bit allocation side yields a fixed-point update whose converged form decomposes into a vessel height, a shared water level, and a finite-alphabet mercury level; the resulting allocation is termed reverse mercury/waterfilling (RMWF).
% Numerical results demonstrate that finite-alphabet awareness in bit allocation sustains end-to-end rate under tight fronthaul budgets and remains robust under antenna scaling — capabilities that become consequential as antenna counts scale faster than fronthaul capacity in modern C-RAN deployments.

The cloud radio access network (C-RAN) mitigates inter-cell interference by jointly processing the observations of distributed remote units (RUs) at a centralized unit (CU), but limited fronthaul capacity forces each RU to compress its received signal. Under transform-compress-forward, an RU transforms its signal and quantizes the resulting coefficients, with bit allocation distributing a finite bit budget across them. Classical reverse waterfilling assumes Gaussian sources, yet practical finite-alphabet symbols carry mutual information that saturates at $\log_2 M$, leaving bit allocation for such inputs unresolved. We address this by formulating bit allocation as maximizing the finite-alphabet generalized mutual information (GMI) achieved after linear MMSE (LMMSE) detection at the CU. Via the I-MMSE relation, this yields a fixed-point update whose converged solution decomposes into a vessel height, a shared water level, and a finite-alphabet mercury level; we term it {reverse mercury/waterfilling} (RMWF). Numerical results show that RMWF sustains end-to-end rate under tight fronthaul budgets and remains robust under antenna scaling, which is increasingly consequential as antenna counts outpace fronthaul capacity in modern C-RAN.
\end{abstract}

\section{Introduction}

Conventional cellular networks process each base station's received signals
independently, treating transmissions from neighboring cells as interference. As networks densify, this inter-cell interference (ICI) becomes a dominant bottleneck. 
The cloud radio access network (C-RAN) architecture~\cite{gesbert2010multi,peng2015fronthaul} addresses this by forwarding observations from distributed remote units (RUs) to a centralized unit (CU), where joint processing mitigates ICI at the network level. 
This cooperative gain, however, is bounded by the fronthaul links between the RUs and the CU: finite capacity forces each RU to compress its observation before forwarding, introducing distortion that degrades performance. 
Efficient fronthaul compression is, therefore, central to realizing the performance gains promised by C-RAN.


From an information-theoretic perspective, C-RAN fronthaul compression has been studied through 
% compress-and-forward and 
distributed source-coding formulations under Gaussian signaling. 
Specifically, in the uplink, the key idea for approaching the network capacity is Wyner--Ziv coding, in which each RU compresses its observation by exploiting the previously decompressed signals at the CU as side information; the CU then performs joint decoding of the user messages over the aggregated reconstructions~\cite{sanderovich2009uplink, park2013robust, zhou2014optimized}. 
% Realizing this scheme requires jointly designed vector-quantized codebooks across RUs, whose encoding complexity grows exponentially with the block length. 
% Setting aside the binning structure that Wyner-Ziv coding demands, even an ordinary vector quantizer at a single RU is already prohibitive for practical fronthaul deployment.
Realizing this scheme requires jointly designed vector-quantized codebooks whose encoding complexity grows exponentially with the block length, making it prohibitive for practical fronthaul, even before the binning structure that Wyner--Ziv coding additionally demands.

% Even an ordinary vector quantizer at a single RU, let alone the binning structure that Wyner--Ziv coding additionally demands, is already prohibitive for real-time fronthaul operation. 
% A more implementation-oriented family of schemes---the one widely adopted in practice---therefore follows the transform-compress-forward paradigm: each RU first applies a transform that maps its received signal to a decorrelated, dimension-reduced representation, then quantizes each component via scalar quantization under the fronthaul budget, and finally forwards the resulting bits to the CU.


% These works characterize what is asymptotically achievable through Wyner--Ziv coding for the uplink and multivariate compression for the downlink, but their achievability relies on high-dimensional vector quantization with codebook complexity growing exponentially in the block length, which is prohibitive for real-time fronthaul links. A more implementation-oriented family of schemes---the one widely adopted in practice---instead follows the transform-compress-forward paradigm: each RU first applies a transform that maps its received signal to a decorrelated, dimension-reduced representation, then quantizes each component via scalar quantization under the fronthaul budget, and finally forwards the resulting bits to the CU.}}

To circumvent the complexity of vector quantization, a widely used practical family of fronthaul compression schemes follows the transform-compress-forward paradigm \cite{liu2015optimized, liu2019twotimescale, zhang2021quantization, qiao2024meta}.
In this paradigm, each RU first applies a transform that maps its received signal to a decorrelated, dimension-reduced representation, then quantizes each coefficient via low-complexity scalar quantization under the fronthaul budget. 
Finally, each quantized representation is forwarded to the CU. 
% each RU first exposes a transformed, decorrelated, or dimension-reduced representation of its received signal to scalar quantizers, and then forwards the resulting bits under the fronthaul budget. 
On the transform side, the Karhunen--Lo\`eve transform (KLT) is a suitable choice, as it decorrelates the received signal along its statistical eigenbasis and concentrates energy on a few dominant coefficients \cite{liu2015optimized, wiffen2020dimension, wiffen2021distributed}.



% The transform side has been extensively studied (Section~\ref{sec:related}); in this work we adopt the standard Karhunen--Lo\`eve transform (KLT) and ask the subsequent question: once the representation is fixed, how should a finite bit budget be allocated across the transformed components?

Once the KLT is adopted as the transform, the remaining crucial design question is how to allocate a finite bit budget across the coefficients. This allocation directly governs how the limited fronthaul capacity is translated into end-to-end rate.
% Information theory provides a clean answer to the sum-distortion minimization problem: 
% for parallel Gaussian sources under a sum rate budget, classical reverse waterfilling~\cite{CoverThomas} assigns more bits to high-variance components, while giving fewer bits to low-variance ones, proportional to the source variances. 
For Gaussian sources, a closely related question of sum-distortion minimization under a total bit budget admits a classical solution via reverse waterfilling (RWF)~\cite{CoverThomas}, which assigns more bits to high-variance coefficients and fewer bits to low-variance ones, according to the logarithm of the source variances.
% in proportion to the source variances. 
This scheme, however, relies on the Gaussian source assumption, which does not hold in modern communication systems. In actual deployments, user data is drawn from finite constellations such as $M$-QAM, so what each RU forwards is a noisy observation of finite-alphabet symbols rather than of a Gaussian source. This makes a clear difference at the bit allocation stage: Gaussian mutual information (MI) grows logarithmically with the signal-to-noise ratio (SNR) without bound, whereas finite-alphabet MI saturates at $\log_2 M$ once the constellation is reliably resolved.



% Assuming Gaussian sources, a classical approach for addressing this is reverse waterfilling~\cite{CoverThomas}, which solves the sum-distortion minimization problem by assigning more bits to high-variance components and fewer bits to low-variance ones, in proportion to the source variances. 
% However, this approach relies on the Gaussian source assumption, which does not hold in modern communication systems. In practice, user data is drawn from finite constellations such as $M$-QAM, so what each RU forwards is a noisy observation of finite-alphabet symbols rather than of a Gaussian source. 
% This makes a clear difference at the bit allocation stage: Gaussian mutual information (MI) grows logarithmically with the signal-to-noise ratio (SNR) without bound, whereas finite-alphabet MI saturates at $\log_2 M$ once the constellation is reliably resolved. A bit allocation scheme designed under the Gaussian assumption does not see this saturation, and ends up spending bits on components that are already resolved while leaving informative ones under-quantized.


% distributes bits according to the source variances, assigning more bits to high-variance components and fewer to low-variance ones. But this answer assumes a Gaussian source, which does not match the practical C-RAN regime: user data is modulated onto finite constellations such as QAM, so what each RU forwards is a noisy and quantized observation of finite-alphabet symbols. The two regimes differ at the bit allocation stage: while Gaussian mutual information (MI) grows logarithmically with the signal-to-noise ratio (SNR), finite-alphabet MI saturates at $\log_2 M$ once the constellation is reliably resolved.A Gaussian-derived rule, blind to this saturation, can over-spends bits on already-resolved components while under-allocating informative ones.

% The same Gaussian-vs-finite-alphabet distinction is in fact the second axis of a four-way duality (Table~\ref{tab:duality}). On the power allocation side, classical waterfilling distributes transmit power across parallel Gaussian channels~\cite{CoverThomas}; for finite-alphabet inputs, mercury/waterfilling modifies this rule by injecting a constellation-dependent mercury level that captures the $\log_2 M$ saturation~\cite{LozanoTulinoVerdu}.
% On the bit allocation side, classical reverse waterfilling addresses the Gaussian case~\cite{CoverThomas}, while the finite-alphabet counterpart---an explicit rule that accounts for the same saturation in the bit budget---has not been developed for bit allocation.
% This missing corner is precisely the operating regime of C-RAN fronthaul compression with QAM signaling.

% The same Gaussian-vs-finite-alphabet distinction is in fact the second axis of a four-way correspondence of MI maximizing allocation rules (Table~\ref{tab:correspondence}).


A parallel Gaussian-vs-finite-alphabet question has, in fact, been resolved on the power allocation side. Classical waterfilling (WF) distributes transmit power across parallel Gaussian channels~\cite{CoverThomas}, and for finite-alphabet inputs, mercury/waterfilling (MWF) modifies this scheme by injecting a constellation-dependent mercury level that captures the finite-alphabet MI saturation~\cite{LozanoTulinoVerdu}. On the bit allocation side, however, a finite-alphabet counterpart that captures the MI saturation remains an open problem, even though this is precisely the operating regime of C-RAN fronthaul compression with finite-alphabet signaling. Motivated by this missing corner, this work proposes a bit allocation scheme that explicitly accounts for finite-alphabet inputs. 

% which motivates the present work: we develop a bit allocation rule that explicitly accounts for the $\log_2 M$ saturation of finite-alphabet inputs under the C-RAN fronthaul compression setting.




% On the bit allocation side, the Gaussian case admits a reverse waterfilling structured allocation---also derivable from classical rate-distortion theory~\cite{CoverThomas}---while the finite-alphabet, MI-aware counterpart, accounting for the same saturation in the bit budget, remains undeveloped. This missing corner is precisely the operating regime of fixed-rate C-RAN fronthaul compression with QAM signaling.

% On the power allocation side, classical waterfilling distributes transmit power across parallel Gaussian channels~\cite{CoverThomas}; for finite-alphabet inputs, mercury/waterfilling modifies this rule by injecting a constellation-dependent mercury level that captures the $\log_2 M$ saturation~\cite{LozanoTulinoVerdu}.
% On the bit allocation side, the Gaussian case admits a reverse waterfilling structured allocation---also derivable from classical rate-distortion theory~\cite{CoverThomas}---while the finite-alphabet, saturation-aware counterpart, accounting for the same saturation in the bit budget, has not been developed for bit allocation.
% On the bit allocation side, the Gaussian case admits a reverse waterfilling structured allocation---also derivable from classical rate-distortion theory~\cite{CoverThomas}---while the finite-alphabet, MI-aware counterpart, accounting for the same saturation in the bit budget, has not been developed for bit allocation.
% This missing corner is precisely the operating regime of fixed-rate C-RAN fronthaul compression with QAM signaling.

% The fourth corner has remained open for two analytical reasons.
% First, finite-alphabet MI and minimum mean squared error (MMSE) admit no closed-form expressions, obstructing the standard Lagrangian--KKT derivation that produced the other three corners.
% Second, in a multi-RU, multi-user C-RAN uplink, the DU's joint receiver couples the effective observations of different users through the bit allocations across multiple RUs, producing a nonconvex problem in which the marginal value of a bit is coupled across users, RUs, and components. 
% % Faced with these obstacles, existing analytical work has understandably relied on Gaussian idealizations or tractable surrogates~\cite{sanderovich2009uplink,park2013robust,park2013joint, zhou2014optimized,liu2015optimized,park2017mixed,rao2015distributed}.
% Faced with these obstacles, existing analytical work has understandably
% favored Gaussian source, where the rate expressions remain tractable. 
% Existing finite-alphabet or waveform-level approaches instead tend to rely on reconstruction-, detection-, or learning-based objectives, yielding useful designs but not an explicit bit-allocation law whose marginal utility is governed by the finite-alphabet post-quantization information gain and its $\log_2 M$ saturation.

% We close the fourth corner directly by taking the post-quantization MI at the DU as the design objective. 
% We close the fourth corner by maximizing the post-linear MMSE (LMMSE) finite-alphabet generalized MI (GMI) at the DU---the achievable end-to-end rate under the standard BICM framework.
% The key analytical lever is the I-MMSE relation~\cite{GuoShamaiVerdu}: although finite-alphabet MI has no closed form, its derivative with respect to SNR equals the finite-alphabet MMSE, which is tabulated for standard constellations.
% This converts the otherwise intractable MI gradient into a tractable marginal-utility expression, yielding a waterfilling-structured fixed-point update---termed reverse mercury/waterfilling (RMWF)---whose converged allocation decomposes into a per-coefficient vessel height, a shared water level, and a finite-alphabet mercury level.
% RMWF extends to multi-RU C-RAN under sum and per-RU fronthaul constraints, recovers the same reverse waterfilling structure in the Gaussian high-SNR limit, and establishes the four-way correspondence summarized in Table~\ref{tab:correspondence}, placing the C-RAN finite-alphabet bit allocation problem on the same analytical footing as its three classical counterparts. 
% The technical contributions are summarized in Section~\ref{sec:contributions}.

\subsection{Related Works}
\label{sec:related}

{\bf{Gaussian Compress-and-Forward:}}
% The earliest C-RAN fronthaul compression studies asked an information-theoretic question: given finite-capacity links between the RUs and the DU, what rates are achievable under Gaussian signaling? 
A foundational line of work characterizes the achievable rates over finite-capacity RU-CU links under Gaussian signaling; a comprehensive treatment of the underlying signal-processing and network-information-theoretic principles is provided in~\cite{park2014fronthaul}.
The uplink direction was first studied through compress-and-forward formulations~\cite{sanderovich2009uplink} and subsequently refined by Wyner--Ziv style robust compression~\cite{park2013robust}, quantization noise covariance design~\cite{zhou2014optimized}, and optimal joint compression-decoding strategies~\cite{zhou2016optimal}. The downlink direction was developed through joint precoding with multivariate compression~\cite{park2013joint} and later through Marton-based generalized compression~\cite{patil2019generalized}, with uplink-downlink duality established in~\cite{liu2021duality}.
% A unified treatment of these results is given in the survey~\cite{park2014fronthaul}.


Building on the same Gaussian abstraction, subsequent works have embedded fronthaul compression into broader network-optimization problems, including joint beamforming and quantization noise covariance design for the uplink~\cite{zhou2016fronthaul}, hybrid data-sharing/compression for the downlink~\cite{patil2018hybrid}, and sparse joint transmission under limited fronthaul capacity~\cite{han2022sparse}. These works treat the bit budget at the abstraction level of test channels, covariance matrices, or fronthaul capacity variables under Gaussian signaling, and an explicit bit allocation scheme for finite-alphabet signals does not arise within this line.



% {\color{red}{A foundational line of work characterizes the rates achievable over finite-capacity RU-CU links under Gaussian signaling, formulated as a network information-theoretic problem.}}
% This was addressed through compress-and-forward and distributed source-coding formulations, beginning with capacity characterizations for the multi-cell  uplink~\cite{sanderovich2009uplink} and extending to Wyner--Ziv-style robust compression~\cite{park2013robust}, quantization noise design for the uplink~\cite{zhou2014optimized}, and joint precoding with multivariate compression for the downlink~\cite{park2013joint}.
% Subsequent work extended this line to optimal uplink compression and decoding strategies~\cite{zhou2016optimal}, uplink--downlink duality~\cite{liu2021duality}, and Marton-based downlink compression~\cite{patil2019generalized}; see~\cite{park2014fronthaul} for a unified survey.

% Building on the same Gaussian abstraction, subsequent work embedded fronthaul compression into broader network-optimization problems, including joint beamforming and quantization noise covariance design for the uplink~\cite{zhou2016fronthaul}, hybrid data-sharing/compression for the downlink~\cite{patil2018hybrid}, and sparse joint transmission under limited fronthaul~\cite{han2022sparse}. 
% These directions operate under the Gaussian signaling assumption and at the abstraction level of test channels, covariance matrices, or fronthaul-rate variables: an explicit bit allocation rule for finite-alphabet signals---the fourth corner of the correspondence in Table~\ref{tab:correspondence}---does not appear within this line.

{\bf{Transform-Compress-Forward and Bit Allocation:}}
A separate line of work makes the transform and scalar-quantization stages explicit, with the bit allocation across transformed coefficients emerging as the central design variable. 
A representative approach is spatial-compression-and-forward for multi-antenna RUs, where each RU applies a linear spatial filter across its antenna array followed by per-coefficient uniform scalar quantization, with the filters designed locally and the transmit powers, receive beamformers, and bit allocation jointly optimized at the CU~\cite{liu2015optimized}. The KLT adopted in this paper can be viewed as a natural realization of such a decorrelating spatial filter.

% Spatial-compression-and-forward applies a linear spatial filter at each RU followed by per-component uniform scalar quantization, with the filters designed locally and the transmit powers, receive beamformers, and bit allocation jointly optimized at the DU~\cite{liu2015optimized}.This line connects naturally to classical decorrelating transforms such as the KLT.

% Related schemes include mixed-analog-to-converter (ADC) architectures that assign different ADC resolutions across antenna chains under a sum resolution constraint guided by low-SNR rate approximation~\cite{park2017mixed}, distributed compressive sensing that reduces RU-side dimension with joint sparse recovery at the DU~\cite{rao2015distributed}, quantization-aware processing that concentrates resources on informative transformed components~\cite{zhang2021quantization}, and two-timescale hybrid compress-and-forward separating long-term analog from short-term digital filtering and bit allocation~\cite{liu2019twotimescale}.

Focusing on the bit allocation aspect, several variants of this paradigm have been investigated: mixed-resolution analog-to-digital converter (ADC) architectures that assign different resolutions across antenna chains under a sum resolution constraint guided by low-SNR rate approximations~\cite{park2017mixed}; distributed compressive sensing that reduces the RU-side dimension and recovers the signal jointly at the CU~\cite{rao2015distributed}; quantization-aware processing that concentrates resources on informative transformed coefficients~\cite{zhang2021quantization}; and two-timescale hybrid compress-and-forward separating long-term analog filtering from short-term digital filtering and bit allocation~\cite{liu2019twotimescale}. However, these works all rely on the Gaussian source assumption and do not capture the saturating MI behavior of finite-alphabet sources.

% Although these methods make the bit budget explicit at the level of dimensions, ADC resolutions, or scalar quantizer bitwidths, 
% % their allocation rules are driven by tractable surrogates---minimum SINR, low-SNR GMI, sparse-recovery guarantees, or long-term Gaussian rate utilities---rather than the finite-alphabet post-quantization MI at the DU, leaving the $\log_2 M$ saturation behavior that defines the correspondence's fourth corner outside their allocation criteria.
% their allocation rules do not yield an explicit information-theoretic marginal utility that captures the $\log_2 M$ saturation defining the correspondence's fourth corner.

{\bf{Learning-Based Transform:}}
In parallel, learning-based approaches have enlarged the design space beyond analytically designed transforms. For instance, deep learning has been adopted as a low-complexity surrogate for joint beamforming and fronthaul quantization design~\cite{yu2021deep}, while progressive distributed compression has been learned from local channel state information (CSI) with a mean squared error (MSE) reconstruction objective~\cite{sohrabi2022learning}.

% In parallel, learning-based approaches have enlarged the design space beyond analytically designed transforms. Specifically, deep learning as a low-complexity surrogate for joint beamforming and fronthaul quantization design~\cite{yu2021deep}, and progressive distributed compression learned from local channel state information (CSI) with an MSE reconstruction objective~\cite{sohrabi2022learning}.

More directly in C-RAN fronthaul compression, meta-learning generates RU-side linear transformation matrices from local CSI, with a gated recurrent unit (GRU)-based refinement stage that uses low-dimensional global feedback to improve sum rate while retaining Gaussian signaling assumptions~\cite{qiao2024meta}.
% A related line---neural source compression---jointly optimizes learned transforms, quantization, and entropy modeling for rate-distortion~\cite{balle2018variational}, and has been extended to neural CPRI compression for modulated waveforms under error vector magnitude (EVM)/rate--distortion rather than communication MI~\cite{bian2025towards}.
A related line is neural source compression, which jointly learns the transform, quantizer, and entropy model under a rate-distortion objective~\cite{balle2018variational}.
This idea has been adapted to neural fronthaul compression for modulated waveforms, adopting an error vector magnitude (EVM)-based distortion measure rather than optimizing communication MI~\cite{bian2025towards}.
% with the objective replaced by error vector magnitude (EVM) and rate-distortion rather than the communication MI~\cite{bian2025towards}.
Closest to the present problem, \cite{chene2024distributed} jointly learns a linear filter, a uniform scalar quantizer with non-homogeneous bit allocation, and a decoder for uplink C-RAN under finite-alphabet signaling and a per-RU bit budget; however, the bit allocation emerges implicitly from black-box supervised training and a bit budget penalty rather than from an interpretable information-theoretic rule.


Across these three lines, the same gap remains: there is no explicit and interpretable fronthaul compression framework for finite-alphabet signaling that captures the saturating MI behavior of finite-alphabet sources.
This paper addresses this gap by directly maximizing the post-quantization MI at the CU under sum or per-RU fronthaul constraints.
Our contribution is therefore not a new transform, quantizer, or entropy model, but an interpretable bit allocation scheme designed specifically to capture this finite-alphabet saturation.
% This paper fills this gap by maximizing the post-quantization MI at the CU under sum or per-RU fronthaul constraints with finite-alphabet signaling.
% Our contribution is therefore not a new transform, quantizer, or entropy model, but an interpretable bit allocation scheme that fills this missing piece.


% Across these three lines, what remains undelivered is the same object: an explicit, interpretable, MI-aware bit allocation rule for finite-alphabet signaling---the missing fourth corner of the correspondence in Table~\ref{tab:correspondence}. 
% This paper supplies that rule by maximizing the post-quantization MI at the DU under sum or per-RU fronthaul constraints with finite-alphabet signaling. 
% The contribution is therefore not a new transform, quantizer, or entropy model, but the interpretable allocation rule that closes the fourth corner and thereby establishes the four-way correspondence.

\subsection{Contributions}
\label{sec:contributions}

% In this paper, we develop a bit allocation scheme that explicitly accounts for finite-alphabet inputs in C-RAN uplink fronthaul compression, filling the missing piece identified above. Our main contributions are as follows.


In this paper, we develop a fronthaul compression method built on the KLT followed by our proposed reverse mercury/waterfilling (RMWF) bit allocation, designed specifically for finite-alphabet inputs in the C-RAN uplink. Our main contributions are as follows.



% We formulate the fronthaul bit allocation problem under
% finite-alphabet inputs as the direct maximization of the end-to-end MI at the DU, evaluated through the GMI of an LMMSE receiver. Two
% operationally distinct fronthaul provisioning regimes---a sum constraint
% modeling pooled provisioning, and a per-RU constraint modeling dedicated
% per-link budgets---are handled within a single algorithmic structure. In
% contrast to prior C-RAN bit allocation works---which operate under the
% Gaussian rate--distortion framework~\cite{liu2015optimized,
% zhang2021quantization}, use low-SNR GMI approximations that miss the
% $\log_2 M$ saturation~\cite{park2017mixed}, or pursue finite-alphabet
% allocation through black-box learned
% surrogates~\cite{chene2024distributed}---the proposed framework adopts the
% finite-alphabet, MI-aware view as the explicit system objective and derives
% the allocation rule analytically through the I-MMSE relation, capturing
% the full finite-alphabet saturation behavior.

% We formulate the fronthaul bit allocation problem under finite-alphabet inputs as the direct maximization of the end-to-end MI at the DU, evaluated through the GMI of an LMMSE receiver. 

{\bf{A General Bit Allocation Framework for Finite-Alphabet C-RAN
Uplinks:}}
We formulate the fronthaul bit allocation problem under finite-alphabet inputs as the maximization of the post-linear minimum mean squared error (LMMSE) finite-alphabet generalized mutual information (GMI)~\eqref{eq:gmi} at the CU.
% , which we adopt as the operational system-level objective.
% ---the achievable end-to-end rate under the standard bit-interleaved coded modulation (BICM) framework, which we adopt as the operational system-level objective.
Two operationally distinct fronthaul provisioning regimes, namely a sum constraint modeling pooled provisioning, and a per-RU constraint modeling dedicated per-link budgets, are handled within a single algorithmic structure. 
% In contrast to prior C-RAN bit allocation works---which operate under the Gaussian rate-distortion framework~\cite{liu2015optimized, zhang2021quantization}, use low-SNR GMI approximations that miss the $\log_2 M$ saturation~\cite{park2017mixed}, or pursue finite-alphabet allocation through black-box learned surrogates~\cite{chene2024distributed}---the proposed framework adopts the finite-alphabet, MI-aware view as the explicit system objective.
Prior C-RAN bit allocation works generally operate under the Gaussian rate-distortion framework~\cite{liu2015optimized, zhang2021quantization}, use low-SNR GMI approximations that miss the $\log_2 M$ saturation~\cite{park2017mixed}, or pursue finite-alphabet allocation through black-box learned surrogates~\cite{chene2024distributed}. 
In contrast, the proposed framework adopts the finite-alphabet, MI-aware view as the explicit system objective.

{\bf{RMWF Structure with a Four-Way Correspondence:}}
% Applying the I-MMSE relation~\cite{GuoShamaiVerdu} on the bit allocation side of the chain, we derive a fixed-point update whose converged allocation decomposes into a per-coefficient vessel height, a shared water level, and a finite-alphabet mercury level---generalizing classical reverse waterfilling through an explicit saturation correction. The allocation reduces exactly to classical reverse waterfilling in the Gaussian limit and, together with mercury/waterfilling for power allocation~\cite{LozanoTulinoVerdu}, completes a four-way correspondence between Gaussian and finite-alphabet inputs on both the power and bit allocation sides of the chain (Table~\ref{tab:correspondence}). Empirically, the iteration converges within 5 to 8 rounds across all tested scenarios, with overall complexity $\mathcal{O}(TK(MN)^3)$.
Applying the I-MMSE relation~\cite{GuoShamaiVerdu} on the bit allocation side of the chain, we derive a fixed-point update whose converged allocation decomposes into a per-coefficient vessel height, a shared water level, and a finite-alphabet mercury level, generalizing the RWF structure through an explicit saturation correction. 
In the Gaussian high-SNR regime, the allocation recovers the same RWF structure (classically derived for rate-distortion~\cite{CoverThomas}); together with WF and MWF on the power allocation side~\cite{LozanoTulinoVerdu}, this establishes a unified MI-maximizing view of the four-way correspondence (Table~\ref{tab:correspondence}), with the finite-alphabet bit allocation corner being the principal new derivation.

{\bf{Numerical Validation under Realistic C-RAN Deployments:}}
The framework is validated on a 3GPP 3D urban-microcell (UMi) C-RAN with hexagonal layouts and mixed modulation orders drawn from $\{4,16,64,256\}$-QAM. Three findings are consistent across all tested configurations: the proposed allocator substantially outperforms both Equal-Bit and Gaussian-input allocations under tight fronthaul budgets; it approaches the unquantized upper bound as the budget grows; and under antenna scaling at fixed fronthaul, it remains on a near-flat plateau while competing schemes degrade. This robustness hierarchy becomes consequential as antenna counts scale faster than fronthaul capacity. The gains carry over from system-level sum GMI to the per-user GMI distribution.
% The iteration converges within 5 to 8 rounds across all tested configurations, with overall complexity $\mathcal{O}(TK(MN)^3)$.

\section{System Model and Problem Formulation}
\label{sec:system_model}

\subsection{C-RAN Uplink with Finite-Alphabet Inputs Model}\label{subsec:cran_fa}

As illustrated in Fig.~\ref{fig:sys}, we consider a C-RAN uplink in which $K$ single-antenna 
users are served by a CU\footnote{We follow the O-RAN terminology \cite{oran2020wp}, where the RU performs RF and low-PHY processing while the CU performs centralized baseband processing. Throughout this paper, we adopt a fully centralized split, under which the high-PHY functions traditionally handled by the distributed unit (DU) are absorbed into the CU.} through $M$ RUs, 
each equipped with $N$ antennas. Each user $k \in 
\CMcal{K} = \{1, \dots, K\}$ is assigned a modulation 
order $M_k \in \{4, 16, 64, 256\}$ corresponding 
to 4-QAM, 16-QAM, 64-QAM, or 256-QAM, and transmits 
with power $P_k$. Let $x_k \in \CMcal{X}_{M_k} 
\subset \mathbb{C}$ denote the symbol transmitted by user $k$, drawn uniformly from the constellation 
$\CMcal{X}_{M_k}$ with 
$\mathbb{E}[|x_k|^2] = P_k$.

The signal observed by the RU $m$ is
\begin{align} \label{eq:received_signal}
  \mathbf{y}_m 
  = \sum_{k=1}^{K} \mathbf{h}_{m,k} x_k 
    + \mathbf{n}_m 
  = \mathbf{H}_m \mathbf{x} + \mathbf{n}_m 
  \in \mathbb{C}^{N},
\end{align}
where $\mathbf{h}_{m,k} \in \mathbb{C}^N$ is the 
channel vector from user $k$ to RU $m$, 
$\mathbf{n}_m \sim \mathcal{CN}(\mathbf{0}, \sigma^2 \mathbf{I}_N)$ is additive Gaussian noise, and the 
aggregated quantities are 
$\mathbf{x} = [x_1, \dots, x_K]^{\sf T} \in 
\mathbb{C}^K$ and $\mathbf{H}_m = [\mathbf{h}_{m,1}, 
\dots, \mathbf{h}_{m,K}] \in \mathbb{C}^{N \times K}$.

Each RU compresses its received signal $\mathbf{y}_m$ into a finite-rate bit stream before forwarding it to the CU.
Following standard practice~\cite{park2014fronthaul}, we adopt a transform coding architecture in which the received signal is first projected onto its statistical eigenbasis via the KLT, and the resulting decorrelated coefficients are then independently quantized. Specifically, let $\mathbf{R}_{\mathbf{y}_m} \triangleq \mathbb{E}[\mathbf{y}_m \mathbf{y}_m^{\sf H}] \in \mathbb{C}^{N \times N}$ denote the received signal covariance matrix at RU $m$, with eigenvalue decomposition
\begin{align}
  \mathbf{R}_{\mathbf{y}_m} = \mathbf{U}_m 
  \boldsymbol{\Lambda}_m \mathbf{U}_m^{\sf H},
\end{align}
where $\mathbf{U}_m \in \mathbb{C}^{N \times N}$ is unitary and $\boldsymbol{\Lambda}_m = \mathrm{diag}(\lambda_{m,1}, \dots, \lambda_{m,N})$ collects the eigenvalues in non-increasing order. 
The KLT coefficients are
\begin{align} \label{eq:klt}
  \boldsymbol{z}_m 
  \triangleq \mathbf{U}_m^{\sf H} \mathbf{y}_m 
  = [z_{m,1}, \dots, z_{m,N}]^{\sf T} 
  \in \mathbb{C}^{N},
\end{align}
which are mutually uncorrelated with $\mathbb{E}[|z_{m,n}|^2] = \lambda_{m,n}$.

% Each coefficient $z_{m,n}$ is then quantized independently with $b_{m,n}$ bits per complex scalar. 
% We adopt the additive distortion model 
% \begin{align} \label{eq:additive_decomp}
%   \hat{z}_{m,n} = z_{m,n} + e_{m,n},
% \end{align}
% where $e_{m,n} \triangleq \hat{z}_{m,n} - z_{m,n}$ is the quantization error. 
% \eqref{eq:additive_decomp} is the standard test-channel form for fronthaul compression and commonly used in C-RAN literature \cite{park2014fronthaul, zhou2014optimized}. 
% Its usefulness for system design relies only on a second-order characterization of $e_{m,n}$. 
% In the high-rate regime, Bennett's white-noise hypothesis~\cite{bennett1948spectra, GershoGray} gives $e_{m,n}$ as zero-mean with variance following the distortion law in~\eqref{eq:qmn_def}. 
% {\color{red}{The way $e_{m,n}$ enters the receiver---and the sense in which its higher-order statistics and its correlation with $z_{m,n}$ are immaterial to the design---is made precise by the mismatched decoding model in Remark~\ref{rem:gaussian_residual}.}}

Each coefficient $z_{m,n}$ is then quantized independently with $b_{m,n}$ bits per complex scalar. We adopt subtractive dithered quantization~\cite{schuchman1964dither}: the encoder at RU~$m$ adds a dither $\delta_{m,n}$ to $z_{m,n}$, applies a uniform scalar quantizer $Q_{m,n}(\cdot)$, and forwards the resulting quantization index to the CU over the fronthaul. 
The CU regenerates the corresponding reproduction level from the index and subtracts the dither realization, yielding
\begin{align}\label{eq:dithered_quant}
  \hat{z}_{m,n} = Q_{m,n}\!\left(z_{m,n} + \delta_{m,n}\right) - \delta_{m,n},
\end{align}
% where $Q_{m,n}(\cdot)$ has cell width $\Delta_{m,n}$, applied separately to the in-phase and quadrature components, and the dither $u_{m,n}$ has real and imaginary parts independently drawn from $\mathrm{Unif}\left([-\Delta_{m,n}/2,\,\Delta_{m,n}/2)\right)$, independent of $\mathbf{x}$ and $\{\mathbf{n}_m\}$.
where $Q_{m,n}(\cdot)$ is a uniform quantizer with cell width $\Delta_{m,n}$ applied separately to the in-phase and quadrature components, and the dither $\delta_{m,n}$ has real and imaginary parts drawn independently from $\mathrm{Unif}\left([-\Delta_{m,n}/2,\,\Delta_{m,n}/2)\right)$, independent of $\mathbf{x}$ and $\{\mathbf{n}_m\}$.
The dither realizations are generated from a pseudorandom seed shared between RU~$m$ and the CU as common randomness, with no per-sample fronthaul overhead.


Defining the quantization error as $e_{m,n} \triangleq \hat{z}_{m,n} - z_{m,n}$, the construction~\eqref{eq:dithered_quant} yields the additive decomposition
\begin{align} \label{eq:additive_decomp}
  \hat{z}_{m,n} = z_{m,n} + e_{m,n}.
\end{align}
Under~\eqref{eq:dithered_quant}, the Schuchman condition~\cite{schuchman1964dither} guarantees that, in the granular (non-overload) regime, $e_{m,n}$ is statistically independent of $z_{m,n}$ and uniformly distributed over $[-\Delta_{m,n}/2,\,\Delta_{m,n}/2)$ on each real dimension, for any source distribution of $z_{m,n}$ and any bit allocation $b_{m,n}$.

Tying the cell width to the coefficient dynamic range via a loading factor $\kappa$, with the per-real-dimension clipping range $A_{m,n} = \kappa\sigma_{r,m,n}$ and $\sigma_{r,m,n}^2 = \lambda_{m,n}/2$ denoting the per-real-dimension variance of $z_{m,n}$, the cell width is
\begin{align} \label{eq:cell_width}
  \Delta_{m,n} = \frac{2 A_{m,n}}{2^{b_{m,n}/2}} = \frac{2\kappa\sigma_{r,m,n}}{2^{b_{m,n}/2}}.
\end{align}
Since the complex error $e_{m,n}$ aggregates the in-phase and quadrature components, each uniform over a cell of width $\Delta_{m,n}$, its second moment admits the following form
\begin{align} \label{eq:qmn_def}
  q_{m,n} 
  \triangleq \mathrm{Var}(e_{m,n})
  = 2\cdot\frac{\Delta_{m,n}^2}{12}
  = c \, \lambda_{m,n} \cdot 2^{-b_{m,n}},
\end{align}
where $c = \kappa^2/3$ is the dithered-uniform quantizer constant; for instance, a $3\sigma$ loading ($\kappa = 3$) yields $c = 3$. The exponential factor $2^{-b_{m,n}}$ reflects the classical distortion-rate scaling of uniform quantization~\cite{bennett1948spectra, gishpierce1968asymptotically, GershoGray}. 
% The high-rate law~\eqref{eq:qmn_def} applies to active coefficients ($b_{m,n} > 0$); coefficients assigned zero bits are dropped and mapped to zero output, in which case $q_{m,n} = \lambda_{m,n}$.

Stacking the quantized coefficients $\hat{\boldsymbol{z}}_m = 
\boldsymbol{z}_m + \mathbf{e}_m$ and applying the inverse transform, the reconstructed signal at the CU is
\begin{align} \label{eq:reconstruction}
  \hat{\mathbf{y}}_m 
  = \mathbf{U}_m \hat{\boldsymbol{z}}_m 
  = \mathbf{y}_m + \mathbf{d}_m,
\end{align}
where $\mathbf{d}_m \triangleq \mathbf{U}_m \mathbf{e}_m$ is the equivalent quantization noise in the original signal domain. Its covariance is 
$\mathbb{E}[\mathbf{d}_m \mathbf{d}_m^{\sf H}] = 
\mathbf{U}_m \mathrm{diag}(q_{m,1}, \dots, q_{m,N}) \mathbf{U}_m^{\sf H}$, which is fully determined by the bit allocation $\{b_{m,n}\}_{n=1}^{N}$ at RU $m$.

\begin{figure*}[t]     
\centerline{\resizebox{1.8\columnwidth}{!}{\includegraphics{Figures/sys.pdf}}}
     \caption{System model of the uplink C-RAN. $K$ users transmit through $\{\mathbf{H}_m\}$ to $M$ RUs with $N$ antennas each. Each RU $m$ projects $\mathbf{y}_m$ via the KLT $\mathbf{U}_m^{\sf H}$, quantizes the coefficients with bits $\{b_{m,n}\}$, and forwards to the CU. The CU applies the inverse transforms $\{\mathbf{U}_m\}$, stacks into $\hat{\mathbf{y}}\in\mathbb{C}^{MN}$, and recovers the $K$ symbols via the LMMSE receiver $\mathbf{W}_{\sf LMMSE}$.}
     \label{fig:sys}
\end{figure*}

\subsection{Joint Decoding Model and Performance Metric}\label{subsec:joint_decoding}
The CU collects the reconstructed signals $\{\hat{\mathbf{y}}_m\}_{m=1}^{M}$ from all RUs and stacks them into the aggregated observation
\begin{align} \label{eq:vmac}
  \hat{\mathbf{y}} 
  \triangleq \big[\, \hat{\mathbf{y}}_1^{\sf T},\, 
       \dots,\, \hat{\mathbf{y}}_M^{\sf T} \,\big]^{\sf T} 
  = \mathbf{H} \mathbf{x} + \mathbf{n} + \mathbf{d} 
  \in \mathbb{C}^{MN},
\end{align}
where $\mathbf{H} = [\mathbf{H}_1^{\sf T}, \dots, \mathbf{H}_M^{\sf T}]^{\sf T} \in \mathbb{C}^{MN \times K}$ is the aggregated channel matrix, $\mathbf{n} = [\mathbf{n}_1^{\sf T}, \dots, \mathbf{n}_M^{\sf T}]^{\sf T} \sim \mathcal{CN}(\mathbf{0}, \sigma^2 \mathbf{I}_{MN})$ is the stacked Gaussian noise, and $\mathbf{d} = [\mathbf{d}_1^{\sf T}, \dots, \mathbf{d}_M^{\sf T}]^{\sf T}$ is the stacked quantization noise with block-diagonal covariance
\begin{align} \label{eq:Q_cov}
  \mathbf{Q} 
  \triangleq \mathbb{E}[\mathbf{d} \mathbf{d}^{\sf H}] 
  = \mathrm{blkdiag}\left(
       \mathbf{U}_1 \mathbf{Q}_1 \mathbf{U}_1^{\sf H},\, 
       \dots,\, 
       \mathbf{U}_M \mathbf{Q}_M \mathbf{U}_M^{\sf H}
     \right),
\end{align}
where $\mathbf{Q}_m \triangleq \mathrm{diag}(q_{m,1}, \dots, q_{m,N})$. Under the dithered construction~\eqref{eq:dithered_quant}, $\mathbf{d}$ is zero-mean and statistically independent of $\mathbf{x}$, so the aggregate disturbance has covariance $\sigma^2\mathbf{I}_{MN}+\mathbf{Q}$.

The CU recovers user symbols by applying a linear receiver to $\hat{\mathbf{y}}$.
Specifically, we adopt the LMMSE receiver given by
\begin{align} \label{eq:lmmse}
  \mathbf{W}_{\sf LMMSE} 
  = \mathbf{P} \mathbf{H}^{\sf H} \mathbf{R}_{\mathbf{y}}^{-1},
\end{align}
where $\mathbf{P} \triangleq \mathrm{diag}(P_1, \dots, P_K)$ and $\mathbf{R}_{\mathbf{y}} \triangleq \mathbf{H} \mathbf{P} \mathbf{H}^{\sf H} + \sigma^2 \mathbf{I}_{MN} + \mathbf{Q}$. 
The soft estimate for user $k$ is then $\tilde{x}_k = [\mathbf{W}_{\sf LMMSE} \hat{\mathbf{y}}]_k$.

The post-equalization SINR for user $k$ admits the standard expression
\begin{align} \label{eq:sinr}
  \mathrm{SINR}_k 
  &= \frac{P_k \, |\mathbf{w}_k^{\sf H} \mathbf{h}_k|^2}
         {\mathbf{w}_k^{\sf H} \big(\sum_{j \neq k} 
            P_j \mathbf{h}_j \mathbf{h}_j^{\sf H} 
            + \sigma^2 \mathbf{I}_{MN} + \mathbf{Q}\big) 
          \mathbf{w}_k},
\end{align}
where $\mathbf{h}_k$ is the $k$-th column of $\mathbf{H}$ and $\mathbf{w}_k^{\sf H}$ is the $k$-th row of $\mathbf{W}_{\sf LMMSE}$. Defining the interference-plus-noise covariance $\mathbf{R}_k \triangleq \sum_{j \neq k} P_j \mathbf{h}_j \mathbf{h}_j^{\sf H} + \sigma^2 \mathbf{I}_{MN} + \mathbf{Q}$ so that $\mathbf{R}_{\mathbf{y}} = \mathbf{R}_k + P_k \mathbf{h}_k\mathbf{h}_k^{\sf H}$, the matrix inversion lemma applied to $\mathbf{R}_{\mathbf{y}}$ in~\eqref{eq:sinr} yields the equivalent compact form $P_k \mathbf{h}_k^{\sf H} \mathbf{R}_k^{-1} \mathbf{h}_k$.

The LMMSE output admits the decomposition $\tilde{x}_k = \rho_k x_k + \eta_k$, where $\rho_k \triangleq \mathbf{w}_k^{\sf H}\mathbf{h}_k$ and $\eta_k$ collects the residual multi-user interference, noise, and quantization distortion. Under the dithered construction, $\eta_k$ is independent of the desired symbol $x_k$, but remains non-Gaussian due to the finite-alphabet multi-user interference.
The decoder applies a Gaussian nearest-neighbor metric, which is mismatched to the true $\eta_k$ and yields the GMI~\cite{merhav1994information, ganti2000mismatched}:
\begin{align} \label{eq:gmi}
  I_k^{\sf GMI} 
  = \sup_{s \ge 0}\;
    \mathbb{E} \!\left[
      \log_2 
      \frac{\exp \!\big(-s\,|\tilde{x}_k - \rho_k x_k|^2\big)}
           {\frac{1}{M_k}\sum_{x' \in \CMcal{X}_{M_k}} 
             \exp \!\big(-s\,|\tilde{x}_k - \rho_k x'|^2\big)}
    \right],
\end{align}
where the expectation is over the joint law of $(x_k, \tilde{x}_k)$ and $M_k$ is the modulation order of user $k$. 
The GMI in~\eqref{eq:gmi} is a lower bound on $I(x_k; \tilde{x}_k)$.

Setting $\gamma_k \triangleq \mathrm{SINR}_k$, we define the constellation-constrained mutual information (CCMI) of a unit-power finite-alphabet input over a scalar AWGN channel~\cite{bruno2026mismatch}:
\begin{align} \label{eq:mi_finite_def}
  I_{M_k}(\gamma_k) 
  &= \log_2 M_k \nonumber \\
  &\quad - \mathbb{E}_{x, n}  \left[
    \log_2 \sum_{x' \in \CMcal{X}_{M_k}} 
    \exp  \left(
      -\frac{|x - x' + n|^2 - |n|^2}{1/\gamma_k}
    \right)
  \right],
\end{align}
with $x$ uniform on $\CMcal{X}_{M_k}$ and $n \sim \mathcal{CN}(0, 1/\gamma_k)$. 
We adopt the CCMI $I_{M_k}(\gamma_k)$ as the system-level design objective throughout this paper; Remark~\ref{rem:gaussian_residual} details its operational meaning.

% \begin{remark}[Operational meaning of the GMI under the dithered residual]
% \label{rem:gaussian_residual}
% \normalfont
% By the mismatched-decoding theorem~\cite{merhav1994information, ganti2000mismatched}, the GMI in~\eqref{eq:gmi} satisfies $I_k^{\sf GMI} \le I(x_k; \tilde{x}_k)$ for any input and noise distribution; this bound is unconditional and requires no specific structure of $\eta_k$. 
% Under the dithered construction, $\eta_k$ is additionally independent of $x_k$.
% In the high-SNR/high-rate regime where RMWF operates, the AWGN test channel is the worst case among additive disturbances of the same covariance, so $I_{M_k}(\gamma_k) \le I_k^{\sf GMI}$ provably---a consequence of the bounded support of the dithered quantization component of the residual (Appendix~\ref{app:bounded_support}).
% For lower rates, a general worst-case-Gaussian guarantee under finite-alphabet inputs is unavailable~\cite{carmon2012disproof}; in this regime, we have observed empirically across all tested configurations that $I_{M_k}(\gamma_k) \le I_k^{\sf GMI}$, consistent with the conservative behavior of Gaussianized decoding metrics reported in~\cite{bruno2026mismatch}. 
% This Gaussian-metric treatment also avoids the exponential complexity of exact joint detection over the true residual law, and mirrors commercial practice: LTE/5G NR receivers compute Gaussian log-likelihood ratios (LLRs) over post-equalization residuals under real-time constraints.
% \end{remark}

% \begin{remark}[Operational meaning of the GMI under the dithered residual]
% \label{rem:gaussian_residual}
% \normalfont
% By the mismatched-decoding theorem~\cite{merhav1994information, ganti2000mismatched}, the GMI in~\eqref{eq:gmi} satisfies $I_k^{\sf GMI} \le I(x_k; \tilde{x}_k)$ for any input and noise distribution; this bound is unconditional and requires no specific structure of $\eta_k$.
% Unfortunately, however, computing $I_k^{\sf GMI}$ directly within the bit-allocation loop is costly, so we instead optimize the CCMI surrogate $I_{M_k}(\gamma_k)$ in~\eqref{eq:mi_finite_def}, which replaces the true residual $\eta_k$ by an independent Gaussian of matched variance.
% Under finite-alphabet inputs a worst-case-Gaussian guarantee does not hold in general, i.e., Gaussianizing the noise need not lower the mutual information~\cite{carmon2012disproof}.
% % and we therefore make no formal claim.
% % The known counterexamples, however, rely on asymmetric or finely structured disturbances, whereas the residual $\eta_k$ here is zero-mean and symmetric, as the multi-user interference, the uniform dither, and the Gaussian noise are each symmetric about the origin. 
% The counterexamples to this conjecture, however, require inputs with large skewness or pronounced excess kurtosis, whereas the residual $\eta_k$ is zero-mean with low-order statistics close to Gaussian---the multi-user interference is a sum of many finite-alphabet terms, and the uniform dither and the Gaussian noise are both symmetric and light-tailed.
% Replacing such a residual by a Gaussian of matched variance is therefore expected to bias the rate estimate conservatively rather than optimistically, consistent with the conservative behavior of Gaussianized decoding metrics observed in~\cite{bruno2026mismatch}. Beyond this conservative tendency, a direct evaluation of the two quantities ($I_k^{\sf GMI}$ and $I_{M_k}(\gamma_k)$) in our C-RAN setting shows them to be virtually indistinguishable across all tested operating points (Section~\ref{subsec:ccmi}). 
% Optimizing the CCMI as the design objective is therefore well justified: it serves as a faithful surrogate whose maximization is empirically equivalent to maximizing the actual GMI in this operating regime. 
% This Gaussian-metric treatment also mirrors commercial practice: LTE/5G NR receivers compute Gaussian log-likelihood ratios (LLRs) over post-equalization residuals under real-time constraints.
% \end{remark}

\begin{remark}[Operational meaning of the GMI under the dithered residual] \label{rem:gaussian_residual}
\normalfont
By the mismatched-decoding theorem~\cite{merhav1994information, ganti2000mismatched},
the GMI in~\eqref{eq:gmi} is an unconditional lower bound on $I(x_k;\tilde{x}_k)$, requiring no structural assumption on $\eta_k$. 
As computing $I_k^{\sf GMI}$ within the bit-allocation loop is costly, we instead optimize the CCMI surrogate $I_{M_k}(\gamma_k)$ in~\eqref{eq:mi_finite_def}, which replaces $\eta_k$ by an independent Gaussian of matched variance. 
Although a worst-case-Gaussian guarantee fails for general finite-alphabet inputs~\cite{carmon2012disproof}, the known counterexamples require large skewness or excess kurtosis.
In our case, $\eta_k$ stays close to Gaussian because it is a sum of many finite-alphabet interferers with symmetric, light-tailed quantization errors induced by subtractive dithering and Gaussian noise. 
Therefore, the surrogate is expected to be conservative~\cite{bruno2026mismatch}. 
A direct evaluation indeed finds $I_k^{\sf GMI}$ and $I_{M_k}(\gamma_k)$ virtually indistinguishable across all
tested points (Section~\ref{subsec:ccmi}), making CCMI maximization empirically equivalent to GMI maximization here. This also mirrors commercial practice: LTE/5G NR receivers
compute Gaussian LLRs over post-equalization residuals under real-time constraints.
\end{remark}

% The GMI thus reflects the rate achievable by the receiver actually deployed in practice, and we adopt it as the system-level objective throughout this paper.

% and matches the operating principle of practical receivers. 
% Moreover, this Gaussian treatment is not merely an analytical convenience: commercial receivers in cellular systems such as LTE and 5G NR routinely compute Gaussian log-likelihood ratios (LLRs) over post-equalization residuals in their modem implementations, since the exact residual law is unavailable and computing one is prohibitive in real time. 

% This Gaussian-metric treatment also avoids the exponential complexity of exact joint detection over the true residual law, and mirrors commercial practice: LTE/5G NR receivers compute Gaussian log-likelihood ratios (LLRs) over post-equalization residuals under real-time constraints. The GMI thus reflects the rate achievable by the receiver actually deployed in practice, and we adopt it as the system-level objective throughout this paper.


% {\color{magenta}\begin{remark}[Rationale behind the Gaussian-and-independent treatment of noise and residuals]
% \label{rem:gaussian_residual}
% \normalfont
% The receiver model adopted above involves two distinct Gaussianity-and-independence assumptions: 
% (i) the quantization noise $\mathbf{q}$ is treated as Gaussian and independent of the data and 
% of the channel noise $\mathbf{n}$ when forming the LMMSE filter, and (ii) the post-equalization 
% residual at the LMMSE output is treated as Gaussian when evaluating the per-user rate through 
% the GMI in~\eqref{eq:gmi}. Neither assumption holds exactly: $\mathbf{q}$ is in general 
% correlated with the source under the pure additive decomposition~\eqref{eq:additive_decomp}, 
% and the post-LMMSE residual is a mixture of finite-alphabet interference, Gaussian noise, and 
% quantization noise. Both assumptions are nonetheless justified on three complementary grounds.

% First, the exact post-LMMSE residual admits no closed-form joint distribution, and its 
% maximum-likelihood decoder requires joint detection with complexity exponential in $K$, which 
% is prohibitive even at moderate scale~\cite{verdu1998multiuser}.

% Second, every deployed cellular receiver treats both the quantization disturbance and the 
% post-equalization residual as Gaussian: this Gaussian-mismatched decoding rule yields simple 
% per-symbol LLR computation and is the de facto interface between the equalizer and the channel 
% decoder in LTE, 5G NR, and forthcoming 6G systems.

% Third, this Gaussian-mismatched decoder admits a rigorous achievable-rate analysis through 
% mismatched-decoding theory. For any second-moment-bounded additive disturbance, Gaussian is 
% the worst-case noise from an information-theoretic standpoint \cite{diggavi2001worst}, so 
% treating $\mathbf{q}$ as Gaussian with covariance $\mathbf{Q}$ can only underestimate the true 
% achievable rate. The corresponding nearest-neighbor decoder, under the mismatched-decoding 
% framework of Lapidoth~\cite{lapidoth1996nearest}, attains the GMI in~\eqref{eq:gmi}, which is 
% always a valid lower bound on $I(\mathbf{x}; \hat{\mathbf{y}})$. The same framework, specialized 
% to finite-alphabet inputs under BICM-type receivers, was developed by \cite{caire1998bicm} and predicts the throughput of practical demodulators.

% Designing the fronthaul compression to maximize~\eqref{eq:gmi} therefore directly optimizes 
% the rate of the receiver actually deployed in the field---not a theoretical upper bound that 
% no real system attains.
% \end{remark}}

\subsection{Problem Formulation}
\label{sec:problem}

% The objective of this work is to determine the bit 
% allocation $\{b_{m,n}\}$---i.e., the number of bits used to quantize the $n$-th KLT coefficient at the $m$-th RU---that maximizes the sum GMI across all $K$ users at the DU. 
% The bit allocation enters the GMI through the quantization noise covariance 
% $\mathbf{Q}$ in~\eqref{eq:Q_cov}, which in turn governs $\mathrm{SINR}_k$ in~\eqref{eq:sinr} and 
% hence $I_k^{\sf GMI}$ in~\eqref{eq:gmi}. We consider two formulations that reflect different fronthaul 
% provisioning regimes encountered in practice.

We seek to find the bit allocation $\{b_{m,n}\}$, i.e., the number of bits quantizing the $n$-th KLT coefficient at RU $m$, for maximizing the sum GMI across all $K$ users at the CU. The allocation enters this objective via $\mathbf{Q}$ in~\eqref{eq:Q_cov}, hence $\mathrm{SINR}_k$~\eqref{eq:sinr} and $I_k^{\sf GMI}$~\eqref{eq:gmi}; we consider two formulations reflecting different fronthaul capacity regimes.

{\textbf{{Sum Fronthaul Bit Allocation:}}}
In some C-RAN deployments, multiple co-located RUs  share a common high-capacity fronthaul link to the CU. 
For example, this occurs when RUs within the same cell site are aggregated through a passive optical network or a shared Ethernet uplink~\cite{checko2015cloud, peng2015fronthaul}. Under such pooled provisioning, the operator may allocate the total bit budget $C_{\sf tot}$ freely across all RUs and all KLT coefficients, which is formulated as
% Under such pooled provisioning, the operator may allocate the total bit budget $C_{\sf tot}$ freely across all RUs and all KLT coefficients.
% The optimal bit allocation solves
% Under such pooled provisioning, the operator may allocate the total bit budget $C_{\sf tot}$ freely across all RUs and all KLT coefficients, which is formulated as 
\begin{align} \label{eq:problem_sum}
  \mathrm{(P1):} \quad 
  \max_{\{b_{m,n}\}} 
  \quad & \sum_{k=1}^{K} 
        I_{M_k} \big(\mathrm{SINR}_k(\{b_{m,n}\})\big) 
        \nonumber \\
  \mathrm{s.t.} 
  \quad & \sum_{m=1}^{M} \sum_{n=1}^{N} b_{m,n} 
        \leq C_{\sf tot}, \\
  & b_{m,n} \geq 0, 
        \quad \forall\, m, n. \nonumber
\end{align}
This setting offers flexibility in bit allocation, since quantization bits can be redistributed both across RUs and across KLT coefficients to maximize the sum GMI.

% This is the most permissive setting and provides an upper bound on what bit allocation alone can achieve in our framework, since the operator may freely 
% shift bits both across RUs (high-quality RUs 
% absorb more bits) and across coefficients
% within each RU (dominant eigenmodes absorb more bits).

{\textbf{{Per-RU Fronthaul Bit Allocation:}}}
In typical deployments, each RU is connected to the 
CU via its own dedicated fronthaul link of capacity 
$C_m$, reflecting independent physical link budgets, service level agreement (SLA) isolation, or distinct fronthaul technologies across RUs.
In this case, the bits at each RU must satisfy a local budget:
\begin{align} \label{eq:problem_per_ru}
  \mathrm{(P2):} \quad 
  \max_{\{b_{m,n}\}} 
  \quad & \sum_{k=1}^{K} 
        I_{M_k} \big(\mathrm{SINR}_k(\{b_{m,n}\})\big) 
        \nonumber \\
  \mathrm{s.t.} 
  \quad & \sum_{n=1}^{N} b_{m,n} \leq C_m, 
        \quad m = 1, \dots, M, \\
  & b_{m,n} \geq 0, 
        \quad \forall\, m, n. \nonumber
\end{align}
Problem~(P2) is structurally harder than~(P1).
% First, the single Lagrange multiplier of~(P1) is replaced by a vector of $M$ multipliers $\{\mu_m\}$, one per RU. 
First, in contrast to~(P1), which involves only a single coupling constraint, (P2) imposes a separate fronthaul capacity constraint for each of the 
$M$ RU-CU links.
Second, bits can no longer be moved across RUs, which removes the main source of flexibility that~(P1) exploits. As a result,~(P2) reduces to $M$ per-RU subproblems that are nonetheless coupled through the joint SINR at the CU, and this coupling is what makes~(P2) difficult to solve.

% Problem (P2) is structurally harder than (P1): the single multiplier of (P1) becomes a vector of $M$ multipliers $\{\mu_m\}$, one per RU, and the cross-RU bit redistribution that (P1) freely exploits is now blocked. Bit allocation under (P2) reduces to $M$ independent per-RU subproblems coupled only through the joint SINR; this coupling is the central source of difficulty in solving (P2).

% Both (P1) and (P2) inherit three structural difficulties from the finite-alphabet setting that distinguish them sharply from their Gaussian-input counterparts. First, $I_M(\cdot)$ admits no closed-form expression and saturates at $\log_2 M$, invalidating the unbounded-growth assumption underlying classical RWF derivations. 
% Second, $\mathrm{SINR}_k$ couples the bit allocations $\{b_{m,n}\}_{m,n}$ across all RUs and all 
% coefficients through the joint linear receiver, so 
% even per-RU allocations cannot be decoupled analytically. Third, the objective is non-convex in 
% $\{b_{m,n}\}$ in general. 
% To address (P1) and (P2) under these challenges, the next section studies an illustrative $N$-parallel channel example that builds the core mechanism of our approach---RMWF.

% Both (P1) and (P2) inherit three finite-alphabet difficulties absent in their Gaussian counterparts:
Both (P1) and (P2) inherit three difficulties:
$I_M(\cdot)$ has no closed form and saturates at $\log_2 M$, breaking the unbounded-growth assumption behind classical RWF; $\mathrm{SINR}_k$ couples $\{b_{m,n}\}$ across all RUs and coefficients through the joint receiver; and the objective is non-convex in general. To address these, the next section studies an illustrative $N$-parallel channel example that builds the core mechanism of our approach, termed RMWF.





\section{Intuitive Example: RMWF} \label{sec:toy_fa}



In this section, we illustrate the key idea behind RMWF through a simple $N$-parallel channel example. 
% We first develop the generalized RMWF solution for practical finite-alphabet inputs such as QAM, and then demonstrate how the classical RWF solution under Gaussian inputs naturally emerges as a special case of this framework.
We first develop the generalized RMWF solution for practical finite-alphabet inputs such as QAM, and then show that for the special case of Gaussian inputs, RMWF reduces to the classical RWF form in the high-SNR regime.

Consider $N$ parallel channels
\begin{align}
  y_i = x_i + n_i, \quad i = 1,\ldots, N
\end{align}
where $n_i \sim \mathcal{CN}(0, \sigma^2)$ is additive Gaussian noise. Each received signal $y_i$ is quantized 
with $b_i$ bits and forwarded to a remote decoder, subject to a total bit budget $\sum_{i=1}^N b_i = B$. 
A relevant question is: how should the bit budget $B$ be distributed across the $N$ channels?


Suppose that each $x_i \in \CMcal{X}_{M_i}$ is drawn from an $M_i$-QAM constellation with $\mathbb{E}[|x_i|^2] = P_i$ for $i = 1, \ldots, N$.
Assuming $b_i$ bits per complex dimension, the quantized observation $\hat{y}_i$ is modeled as $\hat{y}_i = y_i + e_i$ following \eqref{eq:additive_decomp}. 
% where $e_i$ denotes the quantization error, 
% which is 
Also following our quantization model, the quantization error $e_i$ is zero-mean. Its variance is given by
\begin{align} \label{eq:qi}
  q_i \triangleq \mathrm{Var}(e_i) 
  = c \, \sigma_{y_i}^2 \, 2^{-b_i},
\end{align}
where $\sigma_{y_i}^2 = {\rm{Var}}(y_i) = P_i + \sigma^2$. 
% For Gaussian sources, this expression is exact with $c = 1$ under optimal vector quantization (Shannon's rate--distortion bound~\cite{CoverThomas}); for practical scalar quantizers in the high-rate regime, the same form holds with $c > 1$~\cite{GershoGray}. The specific value of $c$ does not affect the qualitative conclusions that follow.
Assuming the decoder uses an MMSE estimator, the resulting post-quantization SNR for estimating $x_i$ is given by $\mathrm{SNR}_i = {P_i}/({\sigma^2 + q_i})$. With the post-quantization SNR, we now turn to the system-level optimization. 
Under a total quantization bit budget $B$, we consider the problem of maximizing the sum GMI between the inputs and the quantized observations:
\begin{align} \label{eq:opt_fa}
  \max_{\{b_i\}} \; \sum_{i=1}^N I_{M_i}(\mathrm{SNR}_i) 
  \quad 
  \text{s.t.} \quad \sum_{i=1}^N b_i \leq B, \; b_i \ge 0, \; \forall i.
\end{align} 
Forming the Lagrangian of~\eqref{eq:opt_fa} with multiplier $\mu > 0$, the KKT necessary condition reads $\partial I_{M_i}(\mathrm{SNR}_i)/\partial b_i \le \mu$ with equality when $b_i > 0$.
By the I-MMSE relation~\cite{GuoShamaiVerdu}, which is applicable under the Gaussian decoding metric of Remark~\ref{rem:gaussian_residual}, the marginal information gain in $\mathrm{SNR}$ admits $dI_{M_i}/d\mathrm{SNR}_i = \mathrm{mmse}_{M_i}(\mathrm{SNR}_i)/\ln 2$, where $\mathrm{mmse}_{M_i}(\cdot)$ is the MMSE function of the $M_i$-QAM constellation. Combined with the chain rule applied to $\mathrm{SNR}_i = P_i/(\sigma^2 + q_i)$, together with $\partial q_i / \partial b_i = -(\ln 2)\, q_i$~\eqref{eq:qi}, the KKT condition reduces to
 % and the chain rule applied to $\mathrm{SNR}_i = P_i/(\sigma^2+q_i)$ with $\partial q_i/\partial b_i = -(\ln 2)\, q_i$, this reduces to the
% exact optimality condition
\begin{align} \label{eq:kkt_fa_exact}
  \frac{P_i \, \mathrm{mmse}_{M_i}(\mathrm{SNR}_i)}{(\sigma^2 + q_i)^2} \cdot q_i 
  \begin{cases}
       = \mu & \text{if } b_i > 0 \\
       \leq \mu & \text{if } b_i = 0
  \end{cases}\,,
\end{align}
where $\mu$ is chosen to satisfy the budget constraint.
Unfortunately, \eqref{eq:kkt_fa_exact} admits no closed-form solution in $b_i$ because $\mathrm{mmse}_{M_i}(\cdot)$ has no analytic expression for finite-alphabet inputs. 
However, fixing the saturation factor $P_i\,\mathrm{mmse}_{M_i}(\mathrm{SNR}_i)/(\sigma^2+q_i)^2$ at the current value reduces~\eqref{eq:kkt_fa_exact} to a residual condition in $q_i$ that admits a closed-form solution. 
This naturally suggests a fixed-point iteration: our RMWF alternates between updating the saturation factor from the current allocation and updating the allocation in closed form, until convergence.

\begin{figure}[t]
    \centering
    \includegraphics[width=0.8\columnwidth]{Figures/LUT/mmse_vs_snr_plot_ver2.eps}
    \caption{
    MMSE function $\mathrm{mmse}_{M_i}(\gamma)$ versus SNR $\gamma$ for representative $M_i \in \{4, 16, 64, 256\}$. 
    }
    \label{fig:mmse_lut}
\end{figure}

Specifically, let $\{b_i^{(t)}\}$ denote the quantization bit allocation at iteration $t$. We evaluate $\mathrm{SNR}_i^{(t)} = P_i/(\sigma^2 + q_i^{(t)})$ 
with $q_i^{(t)} = c (P_i+\sigma^2) \cdot 2^{-b_i^{(t)}}$ and compute the saturation factor\footnote{Since 
$\mathrm{mmse}_{M_i}(\cdot)$ admits no closed-form expression for finite-alphabet inputs, we precompute it over a dense SNR grid for each modulation order and store the values in a lookup table (Fig.~\ref{fig:mmse_lut}), as commonly done in MWF~\cite{LozanoTulinoVerdu}.}
% {\color{red}{XXX}}
\begin{align} \label{eq:mercury_update}
  m_i^{(t)} \triangleq 
  % \mathrm{mmse}_{M_i} \big(\mathrm{SNR}_i^{(t)}\big).
  \frac{P_i \, \mathrm{mmse}_{M_i}(\mathrm{SNR}_i^{(t)})}{(\sigma^2 + q_i^{(t)})^2} .
\end{align}
% Note that the required $\mathrm{mmse}_{M_i}(\mathrm{SNR}_i)$ are evaluated through a precomputed lookup table, shown in Fig.~\ref{fig:mmse_lut}, that tabulates $\mathrm{mmse}_{M_i}(\mathrm{SNR})$ over a dense SNR grid for each modulation order $M_i$.
% Treating $m_i^{(t)}$ as fixed in~\eqref{eq:kkt_fa_exact} 
% reduces the optimality condition to $q_i \cdot m_i^{(t)} = \mu^{(t+1)}$, which yields
Treating $m_i^{(t)}$ as fixed in~\eqref{eq:kkt_fa_exact} reduces the KKT condition to $q_i \cdot m_i^{(t)} = \mu^{(t+1)}$ for $b_i>0$, which yields
\begin{align} \label{eq:bit_update_fa}
b_i^{(t+1)} = 
\max\left\{
0 ,\, 
\log_2{\frac{\lambda_i \cdot m_i^{(t)}}{\mu^{(t+1)}}}
\right\},
\quad
\lambda_i \triangleq c(P_i + \sigma^2).
\end{align}
The multiplier $\mu^{(t+1)}$ is determined by the budget constraint. Let $\CMcal{A}^{(t+1)} = \{i \mid b_i^{(t+1)} > 0\}$ denote the set of active channels at iteration $t+1$. Conditioned on $\CMcal{A}^{(t+1)}$, the multiplier admits the explicit expression
\begin{align} \label{eq:mu_update_fa}
    |\CMcal{A}^{(t+1)}|\log_2 \mu^{(t+1)} 
    = \sum_{i \in \CMcal{A}^{(t+1)}} \log_2 \big( m_i^{(t)} \lambda_i \big) - B.
\end{align}
The complete process of the proposed RMWF is summarized in Algorithm~\ref{alg:rmwf_toy}.



% \begin{align} \label{eq:kkt_fa_frozen}
%   q_i = \frac{\mu^{(t+1)}}{m_i^{(t)}},
% \end{align}
% which has the similar structure as the Gaussian KKT~\eqref{eq:kkt_gauss_exact}. 

% Subsequently, applying the high-rate approximation $(\sigma^2 + q_i)^2 \approx \sigma^4$, \eqref{eq:kkt_fa_frozen} admits the 
% closed-form solution 
% $q_i \approx \mu^{(t+1)} \sigma^4/(P_i m_i^{(t)})$, 
% and substituting into 
% $q_i = c(P_i+\sigma^2) \cdot 2^{-b_i}$ yields the 
% bit update
% \begin{align} \label{eq:bit_update_fa}
%   b_i^{(t+1)} = \log_2 
%   \frac{\lambda_i \cdot m_i^{(t)}}{\mu^{(t+1)}}, 
%   \qquad 
%   \lambda_i \triangleq c P_i(P_i + \sigma^2)/\sigma^4.
% \end{align}
% The multiplier $\mu^{(t+1)}$ is determined by the 
% budget constraint 
% $\sum_{i=1}^N b_i^{(t+1)} = B$, which, owing to 
% the closed-form structure of~\eqref{eq:bit_update_fa}, 
% admits the explicit expression
% \begin{align} \label{eq:mu_update_fa}
%   \log_2 \mu^{(t+1)} 
%   = \tfrac{1}{2} \sum_{i} 
%     \log_2  \big( m_i^{(t)} c(P_i + \sigma^2)\big) 
%     - B/2.
% \end{align}

% The multiplier $\mu^{(t+1)}$ is determined by the budget constraint $\sum_{i=1}^N b_i^{(t+1)} = B$. Let $\CMcal{A}^{(t+1)} = \{i \mid b_i^{(t+1)} > 0\}$ denote the set of active channels at iteration $t+1$. Owing to the closed-form structure of~\eqref{eq:bit_update_fa}, the multiplier admits the explicit expression
% \begin{align} \label{eq:mu_update_fa}
%     \log_2 \mu^{(t+1)} 
%     = \frac{1}{|\CMcal{A}^{(t+1)}|} \sum_{i \in \CMcal{A}^{(t+1)}} \log_2 \big( m_i^{(t)} \lambda_i \big) - \frac{B}{|\CMcal{A}^{(t+1)}|}.
% \end{align}

% % Unlike the Gaussian case~\eqref{eq:b_exact_gauss}, 
% % which required a one-dimensional bisection on $\mu$, 
% In this case, the per-iteration multiplier update here is closed-form. 
% Each iteration thus consists of two elementary steps---\eqref{eq:mercury_update}
% refreshing the first factor from the current allocation, and~\eqref{eq:bit_update_fa}--\eqref{eq:mu_update_fa} updating the allocation in closed form---and the process repeats until the bit allocation $\{b_i^{(t)}\}$ converges. 



% Since $\CMcal{A}^{(t+1)}$ itself depends on $\mu^{(t+1)}$ through~\eqref{eq:bit_update_fa}, the active set and the multiplier must be determined jointly.
% We resolve this by initializing $\CMcal{A}^{(t+1)}$ to the full index set and alternating between~\eqref{eq:mu_update_fa} and~\eqref{eq:bit_update_fa} until $\CMcal{A}^{(t+1)}$ stabilizes; since $\CMcal{A}^{(t+1)}$ evolves over the finite power set $2^{\{1,\dots,N\}}$, this inner loop terminates in finitely many steps.
% % this inner procedure terminates in a finite number of steps because the active set shrinks monotonically. 
Since the active set and the multiplier are interdependent, they must be determined jointly. 
We resolve this by initializing the active set to the full index set and alternating between \eqref{eq:mu_update_fa} and \eqref{eq:bit_update_fa} until the active set stabilizes. 
For a fixed iteration $t$, because removing inactive indices cannot decrease the water level required to satisfy the bit budget, indices discarded within this inner loop cannot be reactivated. 
Thus, the active set can only shrink, and the inner loop terminates after at most $N$ iterations.
Once $\CMcal{A}^{(t+1)}$ is fixed, each iteration of Algorithm~\ref{alg:rmwf_toy} therefore reduces to two elementary steps: \eqref{eq:mercury_update} refreshing the saturation factor $m_i^{(t)}$ from the current allocation, and the joint~\eqref{eq:bit_update_fa}--\eqref{eq:mu_update_fa} update producing the next allocation in closed form. 
This process repeats until $\{b_i^{(t)}\}$ converges.

% Consequently, the multiplier update at each iteration admits a semi-closed-form waterfilling solution. 
% Each iteration thus consists of two elementary steps---\eqref{eq:mercury_update} refreshing the mercury weights from the current allocation, and~\eqref{eq:bit_update_fa}--\eqref{eq:mu_update_fa} jointly determining the active set $\CMcal{A}^{(t+1)}$ and updating the bit allocation the process repeats until the bit allocation $b_i^{(t)}$ converges.

% \begin{algorithm}[t]
% \caption{Reverse Mercury/Waterfilling (RMWF)}
% \label{alg:rmwf_toy}
% \begin{algorithmic}[1]
% \STATE \textbf{Input:} Powers $\{P_i\}$, modulations 
% $\{M_i\}$, noise variance $\sigma^2$, quantization 
% constant $c$, budget $B$, tolerance $\epsilon$
% \STATE \textbf{Initialize:} $b_i^{(0)} \leftarrow 
% B/N$ for all $i$; $\lambda_i \leftarrow c(P_i + \sigma^2)$; 
% $t \leftarrow 0$
% \REPEAT
%   \STATE Compute $q_i^{(t)} = \lambda_i 
%          \cdot 2^{-b_i^{(t)}}$ for all $i$
%   \STATE Compute $\mathrm{SNR}_i^{(t)} = P_i / 
%          (\sigma^2 + q_i^{(t)})$ for all $i$
%   \STATE Update mercury weights $m_i^{(t)}$ 
%          via~\eqref{eq:mercury_update}
%   \STATE Compute $\mu^{(t+1)}$ 
%          via~\eqref{eq:mu_update_fa}
%   \STATE Update $b_i^{(t+1)} $ via~\eqref{eq:bit_update_fa}
%   % \leftarrow 
%   %        \log_2(m_i^{(t)} c(P_i + \sigma^2) / 
%   %               \mu^{(t+1)})$
%   \STATE $t \leftarrow t + 1$
% \UNTIL{$\|\mathbf{b}^{(t)} - \mathbf{b}^{(t-1)}\|_\infty 
%        < \epsilon$}
% \STATE \textbf{Output:} $\{b_i^{(t)}\}$
% \end{algorithmic}
% \end{algorithm}

\begin{algorithm}[t]
\caption{RMWF for $N$-parallel channel}
\label{alg:rmwf_toy}
\begin{algorithmic}[1]
\STATE \textbf{Input:} Powers $\{P_i\}$, modulations 
$\{M_i\}$, noise variance $\sigma^2$, quantization 
constant $c$, budget $B$, tolerance $\epsilon$
\STATE \textbf{Initialize:} $b_i^{(0)} \leftarrow 
B/N$ for all $i$; $\lambda_i \leftarrow c(P_i + \sigma^2)$; 
$t \leftarrow 0$
\REPEAT
  \STATE Compute $q_i^{(t)} = \lambda_i 
         \cdot 2^{-b_i^{(t)}}$ for all $i$
  \STATE Compute $\mathrm{SNR}_i^{(t)} = P_i / 
         (\sigma^2 + q_i^{(t)})$ for all $i$
  \STATE Update $m_i^{(t)}$ 
         via~\eqref{eq:mercury_update}
  \STATE \textbf{Joint $(\mu, \CMcal{A})$ update:} initialize $\CMcal{A}^{(t+1)} \leftarrow \{1, \dots, N\}$
  \REPEAT
    \STATE Compute $\mu^{(t+1)}$ 
           via~\eqref{eq:mu_update_fa} using $\CMcal{A}^{(t+1)}$
    \STATE Update $b_i^{(t+1)}$ via~\eqref{eq:bit_update_fa}
    \STATE $\CMcal{A}^{(t+1)} \leftarrow \{i \mid b_i^{(t+1)} > 0\}$
  \UNTIL{$\CMcal{A}^{(t+1)}$ unchanged}
  \STATE $t \leftarrow t + 1$
\UNTIL{$\|\mathbf{b}^{(t)} - \mathbf{b}^{(t-1)}\|_\infty 
       < \epsilon$}
\STATE \textbf{Output:} $\{b_i^{(t)}\}$
\end{algorithmic}
\end{algorithm}

% At convergence, the iterate $\{b_i^{(t)}\}$ stabilizes at a fixed point $\{b_i^{\star}\}$ with associated mercury weights $m_i^{\star}$ and multiplier $\mu^{\star}$, satisfying the converged form
% \begin{align} \label{eq:rmwf_sol}
%   &b_i^{\star} = \log_2 
%   \frac{\lambda_i \cdot m_i^{\star}}{\mu^{\star}} 
%   = \log_2 \lambda_i
%   + \log_2 m_i^{\star} 
%   - \log_2 \mu^{\star} \\
%   &= \log_2{
%   \left(P_i + \sigma^2\right)} 
%   + \log_2{\left(\frac{P_i}{(\sigma^2+q_i^\star)^2} \cdot \mathrm{mmse}_{M_i}(\gamma_i^\star)
%   \right)} \nonumber \\ 
%   &\qquad  - \log_2 \mu^{\star} \\
%   &= \log_2{\left(\frac{P_i + \sigma^2}{\sigma^2+q_i^\star}\right)} + \log_2{\left(\frac{P_i}{\sigma^2+q_i^\star} \cdot \mathrm{mmse}_{M_i}(\gamma_i^\star)
%   \right)} \nonumber \\
%   &\qquad - \log_2 \mu^{\star} \\
%   &= \log_2{\left(\gamma_i^\star + \frac{\sigma^2}{\sigma^2 + q_i^\star}\right)} 
%   + \log_2{\left(\gamma_i^\star \cdot \mathrm{mmse}_{M_i}(\gamma_i^\star)
%   \right)} \nonumber \\ 
%   &\qquad - \log_2 \mu^{\star} \\
%   &\approx 
%   \underbrace{\log_2{\gamma_i^\star}}_{\text{vessel height}}
%   - \underbrace{\log_2 \mu^{\star}}_{\text{water level}}
%   - \underbrace{\log_2{\frac{1}{\gamma_i^\star \cdot \mathrm{mmse}_{M_i}(\gamma_i^\star)}}}_{\text{mercury level}},
%   % &=
%   % \underbrace{\log_2{\frac{1}{\mu^{\star}}}}_{\text{water level}}
%   % -\underbrace{\log_2{\frac{1}{\gamma_i^\star}}}_{\text{vessel depth}}
%   % - \underbrace{\log_2{\frac{1}{\gamma_i^\star \cdot \mathrm{mmse}_{M_i}(\gamma_i^\star)}}}_{\text{mercury level}}  ,
% \end{align}
% where $\gamma_i^\star = \mathrm{SNR}_i(b_i^\star)$.

% The approximation follows from the high-SNR regime assumption ($\gamma_i^\star \gg 1$). 
% Specifically, since the quantization error is non-negative ($q_i^\star \ge 0$), the fractional term $\frac{\sigma^2}{\sigma^2 + q_i^\star}$ is strictly bounded between $0$ and $1$. 
% Consequently, when the post-quantization SNR is sufficiently high, this term becomes negligible compared to $\gamma_i^\star$.

% This decomposition reveals three ingredients of the allocation. The
% water level $\log_2 \mu^\star$ is a universal threshold set by the
% total budget---bits are assigned only to the excess vessel height above
% it~\cite{CoverThomas}. The vessel height $\log_2\gamma_i^\star$
% captures the post-quantization effective strength, hence depends on the
% final allocation through $q_i^\star$: a higher post-quantization SNR
% yields a taller vessel and more room for bits. The mercury level
% $\log_2(1/(\gamma_i^\star \mathrm{mmse}_{M_i}(\gamma_i^\star)))$ tracks
% the decoder's saturation state: when the constellation is well-resolved,
% $\mathrm{mmse}_{M_i}\to 0$ and the mercury grows, leaving little room for
% additional bits; far from saturation it imposes only a mild correction.

% For a Gaussian source, $\mathrm{mmse}_{\rm G}(\gamma)=1/(1+\gamma)\approx
% 1/\gamma$ at high SNR, so $\gamma\,\mathrm{mmse}_{\rm G}(\gamma)\to 1$
% and the mercury vanishes, recovering $b_i^\star = \log_2(\gamma_i^\star/
% \mu^\star)$---a reverse waterfilling form. Under finite-alphabet inputs,
% the non-vanishing mercury diverts bits from saturated to unsaturated
% channels through the shared water level. Motivated by this analogy with
% mercury/waterfilling for power allocation~\cite{LozanoTulinoVerdu}, we
% term the proposed algorithm reverse mercury/waterfilling (RMWF).

% At convergence, the iterate $\{b_i^{(t)}\}$ stabilizes at a fixed point 
% $\{b_i^{\star}\}$ with associated mercury weights $m_i^{\star}$ and 
% multiplier $\mu^{\star}$. Under the high-SNR regime $\gamma_i^\star \gg 1$, 
% the converged allocation admits the form\footnote{For notational clarity 
% we set $c=1$ in the algebraic decomposition below; the actual algorithm 
% retains $c$ as defined (e.g., $c=\pi\sqrt{3}/2$ for the Panter--Dite 
% scalar quantizer). Reinstating $c$ corresponds to the rescaling 
% $\mu^\star \to c\mu^\star$ and does not affect the structural form.}
% \begin{align} \label{eq:rmwf_sol}
%   &b_i^{\star} = \log_2 
%   \frac{\lambda_i \cdot m_i^{\star}}{\mu^{\star}} 
%   = \log_2 \lambda_i
%   + \log_2 m_i^{\star} 
%   - \log_2 \mu^{\star} \\
%   &= \log_2{
%   \left(P_i + \sigma^2\right)} 
%   + \log_2{\left(\frac{P_i}{(\sigma^2+q_i^\star)^2} \cdot \mathrm{mmse}_{M_i}(\gamma_i^\star)
%   \right)} \nonumber \\ 
%   &\qquad  - \log_2 \mu^{\star} \\
%   &= \log_2{\left(\frac{P_i + \sigma^2}{\sigma^2+q_i^\star}\right)} + \log_2{\left(\frac{P_i}{\sigma^2+q_i^\star} \cdot \mathrm{mmse}_{M_i}(\gamma_i^\star)
%   \right)} \nonumber \\
%   &\qquad - \log_2 \mu^{\star} \\
%   &= \log_2{\left(\gamma_i^\star + \frac{\sigma^2}{\sigma^2 + q_i^\star}\right)} 
%   + \log_2{\left(\gamma_i^\star \cdot \mathrm{mmse}_{M_i}(\gamma_i^\star)
%   \right)} \nonumber \\ 
%   &\qquad - \log_2 \mu^{\star} \\
%   &\approx 
%   \underbrace{\log_2{\gamma_i^\star}}_{\text{vessel height}}
%   - \underbrace{\log_2 \mu^{\star}}_{\text{water level}}
%   - \underbrace{\log_2{\frac{1}{\gamma_i^\star \cdot \mathrm{mmse}_{M_i}(\gamma_i^\star)}}}_{\text{mercury level}},
% \end{align}
% where $\gamma_i^\star = \mathrm{SNR}_i(b_i^\star)$.

% This decomposition reveals three ingredients of the allocation.
% The water level $\log_2 \mu^\star$ is a universal threshold set by the total budget---bits are assigned only to the excess vessel height above it~\cite{CoverThomas}. 
% The vessel height $\log_2 \gamma_i^\star$ captures the post-quantization effective strength, hence depends on the final allocation through $q_i^\star$: a higher post-quantization SNR yields a taller vessel and more room for bits. Notably, unlike mercury/waterfilling for power allocation~\cite{LozanoTulinoVerdu}, where the per-channel gain is fixed by the channel, the vessel height here is self-referential---it depends on the bit allocation $b_i^\star$ itself through $q_i^\star$---a structural feature unique to saturation-aware bit allocation.
% The mercury level $\log_2(1/(\gamma_i^\star \mathrm{mmse}_{M_i}(\gamma_i^\star)))$ tracks the decoder's saturation state: when the constellation is well-resolved, $\mathrm{mmse}_{M_i} \to 0$ and the mercury grows, leaving little room for additional bits; far from saturation it imposes only a mild correction.

% For a Gaussian source, $\mathrm{mmse}_G(\gamma) = 1/(1+\gamma) \approx 1/\gamma$ at high SNR, so $\gamma\,\mathrm{mmse}_G(\gamma) \to 1$ and the mercury vanishes, recovering $b_i^\star = \log_2(\gamma_i^\star/\mu^\star)$---a reverse waterfilling form. 
% Under finite-alphabet inputs, the non-vanishing mercury diverts bits from saturated to unsaturated channels through the shared water level. Motivated by this analogy with mercury/waterfilling for power allocation, we term the proposed algorithm reverse mercury/waterfilling (RMWF).

% The high-SNR approximation in the last line of~\eqref{eq:rmwf_sol} is controlled by the bounded fractional term $\sigma^2/(\sigma^2 + q_i^\star) \in (0,1]$, which is uniformly bounded above by $1$ for any non-negative $q_i^\star$. 
% The condition $\gamma_i^\star \gg 1$ is therefore sufficient for the approximation regardless of the quantization level: coarser quantization (larger $q_i^\star$) only shrinks the fractional term, making the approximation tighter rather than looser. 
% Consequently, the decomposition retains its accuracy in the low-budget regime, where the saturation correction is most consequential.


We now interpret the structure of the proposed RMWF described in Algorithm~\ref{alg:rmwf_toy}.
At convergence, the iterate $\{b_i^{(t)}\}$ stabilizes at a fixed point $\{b_i^{\star}\}$ with associated saturation factor $m_i^{\star}$ and multiplier $\mu^{\star}$. 
Under $\gamma_i^\star \triangleq P_i / (\sigma^2 + q_i^\star) \gg 1$, 
the converged allocation admits the form\footnote{For notational clarity 
we set $c=1$ in the algebraic decomposition below; the actual algorithm 
retains $c$ as defined.
% (e.g., $c=\pi\sqrt{3}/2$ for the Panter-Dite scalar quantizer).
Reinstating $c$ corresponds to the rescaling 
$\mu^\star \to c\mu^\star$ and does not affect the structural form.}
\begin{align} 
  b_i^{\star} 
  % = \log_2 
  % \frac{\lambda_i \cdot m_i^{\star}}{\mu^{\star}} 
  % = \log_2 \lambda_i
  % + \log_2 m_i^{\star} 
  % - \log_2 \mu^{\star} \\
  % &= \log_2{
  % \left(P_i + \sigma^2\right)} 
  % + \log_2{\left(\frac{P_i}{(\sigma^2+q_i^\star)^2} \cdot \mathrm{mmse}_{M_i}(\gamma_i^\star)
  % \right)} \nonumber \\ 
  % &\qquad  - \log_2 \mu^{\star} \\
  &= \log_2{\left(\frac{P_i + \sigma^2}{\sigma^2+q_i^\star}\right)} + \log_2{\left(\frac{P_i}{\sigma^2+q_i^\star} \cdot \mathrm{mmse}_{M_i}(\gamma_i^\star)
  \right)} \nonumber \\
  &\qquad - \log_2 \mu^{\star} \\
  % &= \log_2{\left(\gamma_i^\star + \frac{\sigma^2}{\sigma^2 + q_i^\star}\right)} 
  % + \log_2{\left(\gamma_i^\star \cdot \mathrm{mmse}_{M_i}(\gamma_i^\star)
  % \right)} \nonumber \\ 
  % &\qquad - \log_2 \mu^{\star} \\
  &\approx 
  \underbrace{\log_2{\gamma_i^\star}}_{\text{vessel height}}
  - \underbrace{\log_2 \mu^{\star}}_{\text{water level}}
  - \underbrace{\log_2{\frac{1}{\gamma_i^\star \cdot \mathrm{mmse}_{M_i}(\gamma_i^\star)}}}_{\text{mercury level}}, \label{eq:rmwf_sol}
\end{align}
where the approximation follows from 
\begin{align}
    % \log_2{\left(\gamma_i^\star + \frac{\sigma^2}{\sigma^2 + q_i^\star}\right)} = \log_2 \gamma_i^\star + \log_2{\left( 1+\frac{\sigma^2}{P_i} \right)},
    \log_2{\left(\gamma_i^\star + \frac{\sigma^2}{\sigma^2 + q_i^\star}\right)} 
    \approx \log_2 \gamma_i^\star,
\end{align}
which holds tightly when $\gamma_i^\star$ is sufficiently large, since $q_i \ge 0$.
% which holds tightly when $P_i/\sigma^2$ is sufficiently large.
% Hence the approximation is accurate, when $P_i/\sigma^2$ is sufficiently large.

The above decomposition reveals three crucial ingredients of our allocation.
The water level $\log_2 \mu^\star$ is a universal threshold set by the total budget; bits are assigned only to the excess vessel height above it~\cite{CoverThomas}. 
This plays the same role as the water level in classical RWF, balancing the bit budget across all channels through a single shared threshold.
The vessel height $\log_2 \gamma_i^\star$ captures the post-quantization effective strength, hence depends on the final allocation through $q_i^\star$: a higher post-quantization SNR yields a taller vessel and more room for bits.
% Notably, unlike MWF for power allocation~\cite{LozanoTulinoVerdu}, where the per-channel gain is fixed by the channel, the vessel height here is self-referential. That is to say, it depends on the bit allocation $b_i^\star$ itself through $q_i^\star$. This is a structural feature unique to our RMWF bit allocation. 
Notably, unlike MWF for power allocation~\cite{LozanoTulinoVerdu}, where the per-channel gain is fixed by the channel, the vessel height here is self-referential: it depends on the allocation $b_i^\star$ itself through $q_i^\star$, a structural feature unique to RMWF bit allocation.
The mercury level $\log_2(1/(\gamma_i^\star \mathrm{mmse}_{M_i}(\gamma_i^\star)))$ tracks the decoder's saturation state: when the constellation is well-resolved, $\mathrm{mmse}_{M_i} \to 0$ and the mercury grows, leaving little room for additional bits; far from saturation it imposes only a mild correction.

For a Gaussian source, $\mathrm{mmse}_{\rm G}(\gamma) = 1/(1+\gamma) \approx 1/\gamma$ at high SNR, so $\gamma\,\mathrm{mmse}_{\rm G}(\gamma) \to 1$ and the mercury vanishes. This recovers $b_i^\star = \log_2(\gamma_i^\star/\mu^\star)$, i.e., an RWF form. 
Under finite-alphabet inputs, the non-vanishing mercury diverts bits from saturated to unsaturated channels through the shared water level. Motivated by this analogy with MWF for power allocation, we term the proposed algorithm RMWF.


\begin{remark}[Connection to classical RWF]
\label{rem:classical_rwf}
\normalfont
Classical RWF~\cite{CoverThomas} solves the sum-distortion minimization for $N$ parallel Gaussian sources $\{x_i\}$ with variances $\{\sigma_{x_i}^2\}$ under a total bit budget, yielding $b_i^\star = \max\{0, \log_2(\sigma_{x_i}^2/\mu)\}$. Our formulation maximizes MI rather than minimizing distortion, yet in the Gaussian high-SNR limit it produces the same structural form $b_i^\star = \log_2(\gamma_i^\star/\mu^\star)$.
\end{remark}

\begin{remark}[A four-way correspondence between power and bit allocation]
\normalfont
Together with MWF for power allocation~\cite{LozanoTulinoVerdu}, RMWF completes a four-way correspondence summarized in Table~\ref{tab:correspondence}: the mercury correction emerges on both sides of the communication chain when the Gaussian input assumption is replaced by finite-alphabet constellations.
The Gaussian baseline hides this saturation, since the mercury level vanishes for Gaussian inputs, and the two problems reduce to the classical WF and classical RWF forms, respectively.
% The Gaussian baseline hides this saturation, since $\gamma \cdot \mathrm{mmse}_{\rm G}(\gamma) \rightarrow 1$ at high SNR and the two problems reduce to the classical WF and classical RWF forms, respectively.
% and both problems reduce to their classical WF forms.
\end{remark}

\begin{remark}[Concavity in the high-rate regime]
\label{rem:convexity_toy}
\normalfont
Problem~\eqref{eq:opt_fa} is non-convex in $\{b_i\}$ in general, but concavity is recovered in the high-rate regime. To see this, a direct calculation gives
\begin{align} \label{eq:snr_second_deriv}
    \frac{d^2 \mathrm{SNR}_i}{db_i^2} 
    = (\ln 2)^2 P_i\, q_i \cdot \frac{q_i - \sigma^2}{(\sigma^2 + q_i)^3},
\end{align}
whose sign is governed by $q_i - \sigma^2$. Hence, when $q_i < \sigma^2\, \,\forall i$, $\mathrm{SNR}_i(b_i)$ is concave in each $b_i$. Combined with the concavity and monotonicity of $I_{M_i}(\gamma)$ in $\gamma$~\cite{GuoShamaiVerdu}, the objective in~\eqref{eq:opt_fa} becomes concave in $\{b_i\}$ \cite{boyd2004convex}. 
% The stationary point characterized by~\eqref{eq:kkt_fa_exact} and reached by Algorithm~\ref{alg:rmwf_toy} is therefore a local maximum of~\eqref{eq:opt_fa} whenever the quantization noise stays below the noise floor.
Therefore, the stationary point in the same region is a local maximum.
\end{remark}



% \begin{lemma}[Sufficient concavity condition in the high-rate regime]
% \label{lem:concavity_fa}
% For channel $i$, define
% \begin{align}
%   \gamma_i(b_i)
%   =
%   \frac{P_i}{\sigma^2+q_i(b_i)},
%   \qquad
%   q_i(b_i)
%   =
%   c(P_i+\sigma^2)2^{-b_i}.
% \end{align}
% If
% \begin{align} \label{eq:concavity_threshold_fa}
%   q_i(b_i) < \sigma^2
% \end{align}
% over the operating range of \(b_i\), then the composite function
% \begin{align}
%   b_i \mapsto I_{M_i}(\gamma_i(b_i))
% \end{align}
% is concave in $b_i$.
% Consequently, if
% \eqref{eq:concavity_threshold_fa} holds for all active channels, then the sum objective in~\eqref{eq:opt_fa} is concave over that operating region, and the KKT conditions are sufficient for global optimality within the region.
% \end{lemma}

% \begin{proof}
% By the I-MMSE relation~\cite{GuoShamaiVerdu},
% \begin{align}
%   \frac{d I_M(\gamma)}{d\gamma}
%   =
%   \frac{1}{\ln 2}\mathrm{mmse}_M(\gamma).
% \end{align}
% Since $\mathrm{mmse}_M(\gamma)>0$ and is non-increasing in $\gamma$~\cite{payaro2009hessian}, $I_M(\gamma)$ is concave and non-decreasing in $\gamma$.

% It remains to examine the inner mapping
% \begin{align}
%   \gamma_i(b_i)
%   =
%   \frac{P_i}{\sigma^2+q_i(b_i)},
%   \qquad
%   q_i(b_i)=c(P_i+\sigma^2)2^{-b_i}.
% \end{align}
% Since $dq_i/db_i  =  -(\ln 2)q_i$, we have
% \begin{align}
%   \frac{d\gamma_i}{db_i}
%   &=
%   \frac{(\ln 2)P_iq_i}{(\sigma^2+q_i)^2},
% \end{align}
% and
% \begin{align} \label{eq:d2_snr}
%   \frac{d^2\gamma_i}{db_i^2}
%   =
%   -\frac{
%     P_i(\ln 2)^2 q_i(\sigma^2-q_i)
%   }{
%     (\sigma^2+q_i)^3
%   }.
% \end{align}
% Therefore, if \(q_i<\sigma^2\), then
% \begin{align}
%   \frac{d^2\gamma_i}{db_i^2}<0,
% \end{align}
% so $\gamma_i(b_i)$ is increasing and strictly concave in $b_i$.

% Because $I_M(\gamma)$ is concave and non-decreasing in $\gamma$, the composition rule for concave functions implies that
% $I_{M_i}(\gamma_i(b_i))$ is concave in $b_i$ whenever
% $q_i<\sigma^2$~\cite{boyd2004convex}. 

% If this condition holds for all active channels, then the objective in \eqref{eq:opt_fa} is a sum of concave functions.
% Hence, the KKT conditions are sufficient for global optimality over that region. 
% This completes the proof.
% \end{proof}



\begin{figure}[t]     
\centerline{\resizebox{1.0\columnwidth}{!}{\includegraphics{Figures/Toy/decomposition_snr_centric_matlab_b6.eps}}}
    \caption{Graphical interpretation of RMWF and RWF, based on the high-SNR decomposition in~\eqref{eq:rmwf_sol}. 
    Channel 1 uses 4-QAM with $P_1/\sigma^2 = 20$~dB, and channel 2 uses 256-QAM with $P_2/\sigma^2 = 10$~dB.
    Total quantization bit budget $B=12$.
    The achieved average sum GMI is approximately 5.13 and 4 bits/symbol for RMWF and RWF, respectively, compared to an ideal 5.29 bits/symbol under infinite fronthaul capacity.
    }\label{fig:rmwf_rwf_interp}
\end{figure}

To clearly contrast the RMWF bit allocation incorporating the finite-alphabet inputs against the conventional Gaussian baseline, we illustrate the behaviors of RMWF and RWF in Fig.~\ref{fig:rmwf_rwf_interp}. 
Specifically, we consider a two-channel setup with a total quantization bit budget $B=12$; the first channel is configured with 4-QAM at a high transmit power-to-noise ratio of 20 dB, while the second channel employs 256-QAM at a relatively low 10 dB.
For each channel, the vessel height $\log_2\gamma_i^\star$ encodes the SNR, the mercury level $\log_2(1/(\gamma_i^\star\mathrm{mmse}_{M_i}(\gamma_i^\star)))$ sits above the bottom, and the allocated bits are the clearance between the shared water level $\log_2 \mu^\star$ and the top of the mercury.

The key contrast appears in channel~1. Under RMWF, the high SNR drives $\mathrm{mmse}_4(\gamma_1^\star)\to 0$, producing a thick mercury that blocks additional bit allocation. 
RWF, treating the input as Gaussian, has $\gamma\,\mathrm{mmse}_{\rm G}(\gamma)\to 1$ and therefore zero mercury; it continues pouring bits into the already-saturated channel. 
The vessel heights themselves also differ, because $\gamma_i^\star$ depends on the final allocation: RMWF assigns fewer bits to channel~1 and more bits to channel~2.
For channel 2, the constellation is far from saturation, so the RMWF mercury stays thin and RMWF directs more bits there, whereas RWF, blind to this, does the opposite.
% % the mercury under RMWF is thin, reflected by the non-saturated nature of the 256-QAM constellation; consequently, RMWF redirects more bits to channel~2 than to channel~1. 
% In contrast, RWF fails to account for this saturation characteristic and allocates fewer bits to channel 2.
In short, RMWF discriminates by joint channel strength and modulation-dependent saturation, while RWF discriminates by signal power alone and wastes bits on resolved constellations.

\begin{table*}[t]
\centering
\caption{Four-way correspondence between power and bit allocation for MI maximization.}
\label{tab:correspondence}
\renewcommand{\arraystretch}{1.5}
% \textwidth를 지정하고, \extracolsep{\fill}을 사용해 여백을 자동으로 늘림
\begin{tabular*}{0.85\textwidth}{@{\extracolsep{\fill}} c|cc @{}}
\toprule
 & Gaussian inputs & Finite-alphabet inputs \\
\midrule
Power allocation & Waterfilling & Mercury/waterfilling~\cite{LozanoTulinoVerdu} \\
 & $p_i = \left(1/\mu - 1/g_i\right)^+$ 
 & $p_i = \mathrm{mmse}_{M_i}^{-1}\left(\min\{1, \mu/g_i\}\right) / g_i$ \\
Bit allocation & Reverse waterfilling~\cite{CoverThomas} & \textbf{Reverse mercury/waterfilling (this work)} \\
 & $b_i = \left(\log_2\left(\lambda_i/\mu\right)\right)^+$ 
 % & $b_i = \log_2{1/\mu} - \log_2{1/\gamma_i} - \log_2{1/(\gamma_i\cdot \mathrm{mmse}_{M_i}(\gamma_i))}$ 
 & $b_i = \left(\log_2{\gamma_i} - \log_2{\mu} - \log_2{1/(\gamma_i\cdot \mathrm{mmse}_{M_i}(\gamma_i))}\right)^+$ 
 \\
\bottomrule
\end{tabular*}
\end{table*}




\section{Iterative RMWF-based Fronthaul Compression Framework}
\label{sec:rmwf}

We now extend the RMWF framework developed in the toy example of Section~\ref{sec:toy_fa} to the multi-RU, multi-user C-RAN fronthaul compression problem.
% The following derivation is with respect to the scalar finite-alphabet GMI  surrogate in \eqref{eq:gmi}, not the exact vector mutual information $I(\mathbf{x};\hat{\mathbf{y}})$.
% The I-MMSE relation is applied to the scalar AWGN auxiliary channel parameterized by the post-LMMSE SINR.

\subsection{Sum Fronthaul Constraint}

{\bf{Lagrangian and KKT Condition:}}
Associating the multiplier $\mu > 0$ with the sum fronthaul budget in~\eqref{eq:problem_sum}, the corresponding Lagrangian is given by 
\begin{align} \label{eq:lagrangian_p1}
  \CMcal{L}\big(\{b_{m,n}\}, \mu\big) 
  = \sum_{k=1}^{K} 
    I_{M_k}\big(\mathrm{SINR}_k\big) 
  - \mu  \left(\sum_{m=1}^{M} \sum_{n=1}^{N} 
                b_{m,n} - C_{\sf tot}\right),
\end{align}
% and the stationary condition $\partial \CMcal{L}/\partial b_{m,n} = 0$ requires
and the KKT necessary condition requires
\begin{align} \label{eq:kkt_p1}
  \sum_{k=1}^{K} 
  \frac{d I_{M_k}}{d \mathrm{SINR}_k} 
  \cdot \frac{\partial \mathrm{SINR}_k}
              {\partial b_{m,n}} 
  % = \mu, 
  \begin{cases}
       = \mu & \text{if } b_{m,n} > 0 \\
       \leq \mu & \text{if } b_{m,n} = 0
  \end{cases}\,,
  \quad \forall (m, n).
\end{align}
The first factor in~\eqref{eq:kkt_p1} is obtained by the I-MMSE relation:
\begin{align} \label{eq:immse_user}
  \frac{d I_{M_k}}{d \mathrm{SINR}_k} 
  = \frac{1}{\ln 2} \, 
    \mathrm{mmse}_{M_k}\big(\mathrm{SINR}_k\big).
\end{align}
The second factor requires differentiating the 
post-LMMSE SINR through its dependence on the 
quantization noise covariance $\mathbf{Q}$. 
The following lemma provides the desired expression.

\begin{lemma}[Derivative of post-LMMSE SINR]
\label{lem:dsinr}
Let $\mathrm{SINR}_k = P_k \mathbf{h}_k^{\sf H} 
\mathbf{R}_k^{-1} \mathbf{h}_k$, where $\mathbf{R}_k \triangleq \sum_{j \neq k} P_j \mathbf{h}_j \mathbf{h}_j^{\sf H} + \sigma^2 \mathbf{I}_{MN} + \mathbf{Q}$. 
Then,
\begin{align} \label{eq:dsinr}
  \frac{\partial \mathrm{SINR}_k}{\partial b_{m,n}} 
  = (\ln 2) \, q_{m,n} \cdot P_k \, 
    \big|\mathbf{h}_k^{\sf H} \mathbf{R}_k^{-1} 
         \mathbf{\tilde u}_{m,n}\big|^2,
\end{align}
% where $\mathbf{u}_{m,n} \in \mathbb{C}^{MN}$ 
% denotes the $n$-th eigenvector of RU $m$ embedded 
% in the aggregated coordinate system (i.e., zero 
% outside the $m$-th block).
where $\mathbf{\tilde u}_{m,n} \triangleq \mathbf{v}_m \otimes \mathbf{u}_{m,n} \in \mathbb{C}^{MN}$, with $\mathbf{u}_{m,n} \in \mathbb{C}^N$ being the $n$-th eigenvector of the local covariance matrix $\mathbf{R}_{\mathbf{y}_m}$ at RU $m$, and $\mathbf{v}_m \in \mathbb{R}^M$ denoting the $m$-th standard basis vector.
\end{lemma}

\begin{proof}
The dependence of $\mathrm{SINR}_k$ on $b_{m,n}$ is 
mediated entirely through the quantization noise 
covariance $\mathbf{Q}$. By the chain rule,
\begin{align}
  \frac{\partial \mathrm{SINR}_k}{\partial b_{m,n}} 
  = P_k \mathbf{h}_k^{\sf H} 
    \frac{\partial \mathbf{R}_k^{-1}}
         {\partial b_{m,n}} \mathbf{h}_k.
\end{align}
The matrix derivative identity 
$\partial \mathbf{R}_k^{-1}/\partial b_{m,n} 
= -\mathbf{R}_k^{-1} (\partial \mathbf{R}_k/
\partial b_{m,n}) \mathbf{R}_k^{-1}$ together with 
$\partial \mathbf{R}_k/\partial b_{m,n} 
= \partial \mathbf{Q}/\partial b_{m,n} 
= -(\ln 2) q_{m,n} \mathbf{\tilde u}_{m,n} \mathbf{\tilde u}_{m,n}^{\sf H}$, where the last equality follows from $q_{m,n} = c \lambda_{m,n} \cdot 2^{-b_{m,n}}$ and the block structure of $\mathbf{Q}$ gives
\begin{align}
  \frac{\partial \mathbf{R}_k^{-1}}{\partial b_{m,n}} 
  = (\ln 2) q_{m,n} \mathbf{R}_k^{-1} 
    \mathbf{\tilde u}_{m,n} \mathbf{\tilde u}_{m,n}^{\sf H} 
    \mathbf{R}_k^{-1}.
\end{align}
Substituting this into the chain rule expression 
yields~\eqref{eq:dsinr}.
\end{proof}

Substituting~\eqref{eq:immse_user} 
and~\eqref{eq:dsinr} into~\eqref{eq:kkt_p1} yields the 
exact KKT condition
\begin{align} \label{eq:kkt_p1_exact}
  q_{m,n} \cdot S_{m,n}\big(\{b_{m,n}\}\big) 
  % = \mu, 
  \begin{cases}
       = \mu & \text{if } b_{m,n} > 0 \\
       \leq \mu & \text{if } b_{m,n} = 0
  \end{cases}\,,
  \quad \forall (m, n),
\end{align}
where $S_{m,n}\big(\{b_{m,n}\}\big) $ is the user-aggregated saturation score defined as
\begin{align} \label{eq:score_def}
  S_{m,n}\big(\{b_{m,n}\}\big) 
  \triangleq \sum_{k=1}^{K} 
    P_k \, \mathrm{mmse}_{M_k}\big(\mathrm{SINR}_k\big) 
    \cdot \big|\mathbf{h}_k^{\sf H} 
                \mathbf{R}_k^{-1} 
                \mathbf{\tilde u}_{m,n}\big|^2.
\end{align}
% The score $S_{m,n}$ generalizes the toy-example mercury weight $m_i$~\eqref{eq:mercury_update} to the multi-user, multi-RU C-RAN setting. Each user $k$ contributes its $\mathrm{mmse}_{M_k}(\mathrm{SINR}_k)$ scaled by the LMMSE projection $|\mathbf{h}_k^{\sf H} \mathbf{R}_k^{-1} 
% \mathbf{\tilde u}_{m,n}|^2$, which quantifies how strongly the $(m,n)$-th KLT coefficient influences user $k$'s detection. 
% Concretely, the score weighs each user's saturation by how much the $(m,n)$-th coefficient matters to that user's detection: users that heavily rely on this coefficient contribute their full mercury weight, while users that barely use it contribute almost nothing.
The score $S_{m,n}$ generalizes the toy-example mercury weight $m_i$~\eqref{eq:mercury_update}
to the multi-user, multi-RU C-RAN setting.
The $(m,n)$-th coefficient contributes to each user $k$ through the product of two factors: the LMMSE projection $|\mathbf{h}_k^{\sf H}\mathbf{R}_k^{-1}\mathbf{\tilde u}_{m,n}|^2$, measuring how strongly user $k$'s detection leverages it, and $P_k\,\mathrm{mmse}_{M_k}(\mathrm{SINR}_k)$, measuring how much a marginal SINR gain is still worth under user $k$'s saturation state.
Their product is the coefficient's marginal contribution to user $k$'s GMI, so $S_{m,n}$ scores the total contribution of the $(m,n)$-th coefficient to the sum GMI.





{\bf{Iterative RMWF Algorithm:}}
The KKT condition~\eqref{eq:kkt_p1_exact} couples 
the bit allocations $\{b_{m,n}\}$ through the score 
$S_{m,n}$, which itself depends on $\{b_{m,n}\}$ via 
the post-LMMSE SINRs $\{\mathrm{SINR}_k\}$ (through 
$\mathbf{R}_k$). 
% This coupling is more involved 
% than in the toy example, where each $m_i$ depended 
% only on $b_i$, but the same structural separation 
% applies: the score $S_{m,n}$ plays the role of the 
% mercury weight, while the quantization noise factor 
% $q_{m,n}$ admits a closed-form inversion when the 
% score is held fixed. We therefore adopt the same 
% fixed-point philosophy of Section~\ref{sec:toy_fa}, 
% alternating between score evaluation and bit-update 
% steps.
Specifically, given the iterate $\{b_{m,n}^{(t)}\}$ at step $t$, we compute the post-quantization SINRs $\{\mathrm{SINR}_k^{(t)}\}$ via~\eqref{eq:sinr} with $\mathbf{Q}^{(t)}$ assembled from $q_{m,n}^{(t)} = c \lambda_{m,n} \cdot 2^{-b_{m,n}^{(t)}}$, and refresh the scores
\begin{align} \label{eq:score_update}
  S_{m,n}^{(t)} 
  \triangleq \sum_{k=1}^{K} 
    P_k \, \mathrm{mmse}_{M_k}\big(\mathrm{SINR}_k^{(t)}\big) 
    \cdot \big|\mathbf{h}_k^{\sf H} 
              \big(\mathbf{R}_k^{(t)}\big)^{-1} 
              \mathbf{\tilde u}_{m,n}\big|^2.
\end{align}
By treating $\{S_{m,n}^{(t)}\}$ as fixed in~\eqref{eq:kkt_p1_exact}, the KKT condition reduces to
\begin{align} \label{eq:bit_update_p1}
  b_{m,n}^{(t+1)} 
  = \max\left\{
  0,\,
  \log_2 \frac{c\lambda_{m,n} \cdot S_{m,n}^{(t)}}
                {\mu^{(t+1)}}\right\},
\end{align}
where we have substituted $q_{m,n} = c \lambda_{m,n} \cdot 2^{-b_{m,n}}$ and 
written $\lambda_{m,n}$ for the KLT eigenvalue. 
% Comparing~\eqref{eq:bit_update_p1} with the toy-example bit update~\eqref{eq:bit_update_fa} reveals an identical structure: 
% {\color{magenta}
% Specifically, the per-channel allocation is a vessel depth $\log_2 \lambda_{m,n}$ augmented by a saturation correction $\log_2 S_{m,n}^{(t)}$ and offset by a shared water level $\log_2 \mu^{(t+1)}$.}

% The multiplier $\mu^{(t+1)}$ is determined by 
% imposing the sum-rate budget 
% $\sum_{m,n} b_{m,n}^{(t+1)} = C_{\sf tot}$ 
% on~\eqref{eq:bit_update_p1}, which admits the 
% explicit expression
% \begin{align} \label{eq:mu_closed}
%   \log_2 \mu^{(t+1)} 
%   = \frac{1}{MN} 
%     \sum_{m=1}^{M} \sum_{n=1}^{N} 
%     \log_2 \big(\lambda_{m,n} \cdot S_{m,n}^{(t)}\big) 
%   - \frac{C_{\sf tot}}{MN}.
% \end{align}
% As in the toy example, no bisection on $\mu$ is 
% required. Each iteration thus consists of two 
% elementary steps---\eqref{eq:score_update} refreshing the scores, 
% and~\eqref{eq:bit_update_p1}--\eqref{eq:mu_closed} 
% updating the allocation in closed form---and the 
% process repeats until convergence. The full 
% procedure is summarized in 
% Algorithm~\ref{alg:rmwf_p1}.

The multiplier $\mu^{(t+1)}$ is determined by 
imposing the sum-rate budget 
$\sum_{m,n} b_{m,n}^{(t+1)} = C_{\sf tot}$ 
on~\eqref{eq:bit_update_p1}. Let 
$\CMcal{A}^{(t+1)} = \{(m, n) \mid b_{m,n}^{(t+1)} > 0\}$ 
denote the set of active KLT coefficients at 
iteration $t+1$. Conditioned on 
$\CMcal{A}^{(t+1)}$, the multiplier admits the 
explicit expression
\begin{align} \label{eq:mu_closed}
  % \log_2 \mu^{(t+1)} 
  % = \frac{1}{|\CMcal{A}^{(t+1)}|} 
  %   \sum_{(m,n) \in \CMcal{A}^{(t+1)}} 
  %   \log_2 \big(\lambda_{m,n} \cdot S_{m,n}^{(t)}\big) 
  % - \frac{C_{\sf tot}}{|\CMcal{A}^{(t+1)}|}.
  |\CMcal{A}^{(t+1)}|\log_2 \mu^{(t+1)} 
  = \sum_{(m,n) \in \CMcal{A}^{(t+1)}} 
    \log_2 \big(c\lambda_{m,n} \cdot S_{m,n}^{(t)}\big) 
  - C_{\sf tot}.
\end{align}
Since $\CMcal{A}^{(t+1)}$ itself depends on $\mu^{(t+1)}$ through~\eqref{eq:bit_update_p1}, the active set and the multiplier must be determined jointly. 
We do this by initializing $\CMcal{A}^{(t+1)}$ to the full index set $\{(m, n)\}_{m,n}$ and alternating  between~\eqref{eq:mu_closed} and~\eqref{eq:bit_update_p1} until $\CMcal{A}^{(t+1)}$ stabilizes. 
As in the toy example, this inner loop terminates finitely and no bisection on $\mu$ is required. 
Each iteration thus consists of two elementary steps: \eqref{eq:score_update} refreshing the scores, and~\eqref{eq:bit_update_p1}--\eqref{eq:mu_closed} updating the allocation in closed form.
This process repeats until convergence. The full procedure is summarized in Algorithm~\ref{alg:rmwf_p1}.



The per-iteration complexity is dominated by the $K$ matrix inverses 
$(\mathbf{R}_k^{(t)})^{-1} \in \mathbb{C}^{MN \times MN}$ in the SINR step, 
costing $\mathcal{O}(K(MN)^3)$. Empirically, the iteration converges in 
$T = 5$--$8$ steps for typical C-RAN configurations 
(see Section~\ref{subsec:convergence}), making the overall complexity 
$\mathcal{O}(T K (MN)^3)$. 
% As in the toy example~(Section~\ref{sec:toy_fa}), formal convergence guarantees are not established; empirical convergence is consistently observed across all tested scenarios.



% \begin{remark}[Convergence]
% \label{rem:convergence}
% \normalfont
% The convergence of Algorithm~\ref{alg:rmwf_p1} is supported by two structural observations. The closed-form multiplier update in~\eqref{eq:mu_closed} keeps every iterate feasible, and the bit allocation is bounded in $[0, C_{\sf tot}]$, confining $\{\mathbf{b}^{(t)}\}$ to a compact set; the Bolzano-Weierstrass theorem therefore guarantees the existence of a convergent subsequence. Convergence to the same fixed point within $T=5$--$8$ iterations is empirically verified in Section~\ref{subsec:convergence}.
% \end{remark}

\begin{algorithm}[t]
\caption{Iterative RMWF for (P1)}
\label{alg:rmwf_p1}
\begin{algorithmic}[1]
\STATE \textbf{Input:} Channels $\{\mathbf{H}_m\}$, 
powers $\{P_k\}$, modulations $\{M_k\}$, 
eigenvalues $\{\lambda_{m,n}\}$, eigenvectors 
$\{\mathbf{\tilde u}_{m,n}\}$, budget $C_{\sf tot}$, 
tolerance $\epsilon$
\STATE \textbf{Initialize:} $b_{m,n}^{(0)} 
\leftarrow C_{\sf tot}/(MN)$ for all $(m, n)$; 
$t \leftarrow 0$
\REPEAT
  \STATE Form $\mathbf{Q}^{(t)} = 
         \mathrm{blkdiag}(\mathbf{U}_m \mathbf{Q}_m^{(t)} 
         \mathbf{U}_m^{\sf H})$ with 
         $q_{m,n}^{(t)} = c \lambda_{m,n} \cdot 
         2^{-b_{m,n}^{(t)}}$
  \STATE Compute $\mathrm{SINR}_k^{(t)} 
         = P_k \mathbf{h}_k^{\sf H} 
         (\mathbf{R}_k^{(t)})^{-1} \mathbf{h}_k$ 
         for all $k$
  \STATE Update scores $S_{m,n}^{(t)}$ 
         via~\eqref{eq:score_update}
  \STATE \textbf{Joint $(\mu, \CMcal{A})$ update:} 
         initialize $\CMcal{A}^{(t+1)} \leftarrow 
         \{(m, n) \mid 1 \le m \le M, 1 \le n \le N\}$
  \REPEAT
    \STATE Compute $\mu^{(t+1)}$ 
           via~\eqref{eq:mu_closed} using 
           $\CMcal{A}^{(t+1)}$
    \STATE Update $b_{m,n}^{(t+1)}$ 
           via~\eqref{eq:bit_update_p1}
    \STATE $\CMcal{A}^{(t+1)} \leftarrow 
           \{(m, n) \mid b_{m,n}^{(t+1)} > 0\}$
  \UNTIL{$\CMcal{A}^{(t+1)}$ unchanged}
  \STATE $t \leftarrow t + 1$
\UNTIL{$\|\mathbf{b}^{(t)} - 
       \mathbf{b}^{(t-1)}\|_{\infty} < \epsilon$}
% \STATE Project $\{b_{m,n}\}$ onto non-negative 
%        integers (if required)
\STATE \textbf{Output:} $\{b_{m,n}^{(t)}\}$
\end{algorithmic}
\end{algorithm}


% \begin{remark}[Concavity of (P1)]
% \label{rem:convexity_p1}
% \normalfont
% The toy-example concavity 
% result~(Lemma~\ref{lem:concavity_fa}) extends to 
% problem~\eqref{eq:problem_sum} via the same 
% composition argument applied per coefficient: when 
% $q_{m,n} < \sigma^2$ for all $(m,n)$, the per-user 
% GMI $I_{M_k}(\mathrm{SINR}_k(\{b_{m,n}\}))$ is 
% concave in $\{b_{m,n}\}$, and the sum objective is 
% therefore concave. Under this high-rate condition, 
% the KKT condition~\eqref{eq:kkt_p1_exact} is 
% sufficient for global optimality and 
% Algorithm~\ref{alg:rmwf_p1} converges to the global 
% maximum. This regime covers the practical 
% fronthaul operating point---per-coefficient 
% budgets of several bits at moderate per-coefficient 
% SNRs---under which the proposed RMWF solution is 
% globally optimal.
% \end{remark}


\subsection{Per-RU Fronthaul Constraint}



Problem~(P2) inherits the KKT structure and fixed-point iteration of
Section~\ref{sec:rmwf}-A verbatim; only the global multiplier $\mu$
splits into $M$ RU-local multipliers $\{\mu_m\}$. Associating $\mu_m > 0$ with the $m$-th constraint in~\eqref{eq:problem_per_ru} yields the
localized KKT condition
\begin{align} \label{eq:kkt_p2_exact}
  q_{m,n} \cdot S_{m,n}\big(\{b_{m,n}\}\big)
  \begin{cases}
       = \mu_m & \text{if } b_{m,n} > 0 \\
       \leq \mu_m & \text{if } b_{m,n} = 0
  \end{cases},
  \quad \forall (m, n),
\end{align}
with $S_{m,n}$ defined as in~\eqref{eq:score_def}. The bits are then
updated by
\begin{align} \label{eq:bit_update_p2}
  b_{m,n}^{(t+1)} = \max \left\{0,\,
    \log_2 \frac{c\lambda_{m,n} \cdot S_{m,n}^{(t)}}
                {\mu_m^{(t+1)}}\right\},
\end{align}
where each $\mu_m^{(t+1)}$ is determined by imposing the local budget
$\sum_n b_{m,n}^{(t+1)} = C_m$, yielding the closed form
\begin{align} \label{eq:mu_closed_p2}
  |\CMcal{A}_m^{(t+1)}| \log_2 \mu_m^{(t+1)}
  = \sum_{n \in \CMcal{A}_m^{(t+1)}}
    \log_2 \big(c\lambda_{m,n} \cdot S_{m,n}^{(t)}\big) - C_m,
\end{align}
with $\CMcal{A}_m^{(t+1)}=\{n\mid b_{m,n}^{(t+1)}>0\}$. As in~(P1),
each $(\mu_m, \CMcal{A}_m)$ pair is determined jointly by alternation;
the inner loop now runs independently per RU. Structurally,~(P2) decomposes into $M$ local RMWF problems coupled indirectly through the scores: a bit-starved user inflates the scores of all coefficients serving it across all RUs, prompting each RU to reallocate its local budget toward that user. The full procedure appears in Algorithm~\ref{alg:rmwf_p2}, with overall complexity $\mathcal{O}(TK(MN)^3)$ identical to~(P1).

% {\color{red}{
% \begin{remark}[Concavity of (P1) and (P2)]
% \label{rem:convergence}
% \normalfont
% The convergence of Algorithm~\ref{alg:rmwf_p1} is supported by three structural observations. The closed-form multiplier update in~\eqref{eq:mu_closed} keeps every iterate feasible; the bit allocation is bounded in $[0, C_{\sf tot}]$, confining $\{b_{m,n}^{(t)}\}$ to a compact set; and in the high-rate regime $q_{m,n} \ll \sigma^2$, the composition-based concavity argument of Remark~\ref{rem:convexity_toy} extends to the multi-user, multi-RU setting, rendering the objective concave on the relevant region. Convergence to the same fixed point within $T=5$--$8$ iterations is verified empirically in Section~\ref{subsec:convergence}.
% \end{remark}
% }}

% \begin{remark}[Concavity of (P1) and (P2)]
% \label{rem:concavity_cran}
% \normalfont
% Problems (P1) and (P2) are non-convex in $\{b_{m,n}\}$ in general, 
% % because each $b_{m,n}$ couples all users' $\mathrm{SINR}_k$ through $\mathbf{R}_k$, 
% but concavity is recovered in the high-rate regime $q_{m,n} < \sigma^2\,\, \forall m,n$. 
% In this regime, the Loewner bound 
% $\mathbf{R}_k^{-1} \preceq (\sigma^2 \mathbf{I}_{MN} + \mathbf{Q})^{-1}$ 
% together with the KLT-diagonal structure of $\mathbf{Q}$ yields 
% $\nabla^2 \mathrm{SINR}_k \preceq \mathbf{0}$; combined with the 
% monotonicity and concavity of $I_{M_k}(\gamma)$~\cite{GuoShamaiVerdu}, 
% the objective Hessian is negative semidefinite, so the KKT fixed point 
% reached by Algorithm~\ref{alg:rmwf_p1} and~\ref{alg:rmwf_p2} coincides with the local
% optimum.
% \end{remark}

\begin{remark}[Concavity of (P1) and (P2)]
\label{rem:concavity_cran}
\normalfont
Problems (P1) and (P2) are non-convex in $\{b_{m,n}\}$ in general, 
but concavity is recovered in the high-rate regime $q_{m,n} < \sigma^2\,\, \forall m,n$. 
To establish $\nabla^2 \mathrm{SINR}_k \preceq \mathbf{0}$, 
$\mathbf{R}_k^{-1} \preceq (\sigma^2 \mathbf{I}_{MN} + \mathbf{Q})^{-1}$ 
together with the KLT-diagonal structure of $\mathbf{Q}$ reduces the 
second directional derivative of $\mathrm{SINR}_k$ to the 
per-coefficient comparison 
$2q_{m,n}^2/(\sigma^2+q_{m,n}) \le q_{m,n}$, equivalent to 
$q_{m,n} \le \sigma^2$. 
The subsequent composition with the concavity and monotonicity of $I_{M_k}$ then proceeds identically to Remark~\ref{rem:convexity_toy}, so the stationary point in the same region is a local maximum.
\end{remark}

\begin{algorithm}[t]
\caption{Iterative RMWF for (P2)}
\label{alg:rmwf_p2}
\begin{algorithmic}[1]
\STATE \textbf{Input:} Channels $\{\mathbf{H}_m\}$, 
powers $\{P_k\}$, modulations $\{M_k\}$, 
eigenvalues $\{\lambda_{m,n}\}$, eigenvectors 
$\{\mathbf{\tilde u}_{m,n}\}$, per-RU budgets $\{C_m\}$, 
tolerance $\epsilon$
\STATE \textbf{Initialize:} $b_{m,n}^{(0)} 
\leftarrow C_m/N$ for all $(m, n)$; $t \leftarrow 0$
\REPEAT
  \STATE Form $\mathbf{Q}^{(t)} = 
         \mathrm{blkdiag}(\mathbf{U}_m \mathbf{Q}_m^{(t)} 
         \mathbf{U}_m^{\sf H})$ with 
         $q_{m,n}^{(t)} = c \lambda_{m,n} \cdot 
         2^{-b_{m,n}^{(t)}}$
  \STATE Compute $\mathrm{SINR}_k^{(t)}$ 
         and update scores $S_{m,n}^{(t)}$ 
         via~\eqref{eq:score_update}
  \FOR{$m = 1, \dots, M$}
    \STATE \textbf{Joint $(\mu_m, \CMcal{A}_m)$ update:} 
           initialize $\CMcal{A}_m^{(t+1)} \leftarrow 
           \{1, \dots, N\}$
    \REPEAT
      \STATE Compute $\mu_m^{(t+1)}$ 
             via~\eqref{eq:mu_closed_p2} using 
             $\CMcal{A}_m^{(t+1)}$
      \STATE Update $b_{m,n}^{(t+1)}$ 
             via~\eqref{eq:bit_update_p2} for all $n$
      \STATE $\CMcal{A}_m^{(t+1)} \leftarrow 
             \{n \mid b_{m,n}^{(t+1)} > 0\}$
    \UNTIL{$\CMcal{A}_m^{(t+1)}$ unchanged}
  \ENDFOR
  \STATE $t \leftarrow t + 1$
\UNTIL{$\|\mathbf{b}^{(t)} - 
       \mathbf{b}^{(t-1)}\|_{\infty} < \epsilon$}
% \STATE Project $\{b_{m,n}\}$ onto non-negative integers 
%        (if required)
\STATE \textbf{Output:} $\{b_{m,n}^{(t)}\}$
\end{algorithmic}
\end{algorithm}


% \begin{remark}[Concavity of (P2)]
% \label{rem:convexity_p2}
% \normalfont
% The concavity argument of Remark~\ref{rem:convexity_p1} applies 
% verbatim: when $q_{m,n} < \sigma^2$ for all $(m,n)$, the sum-GMI 
% objective is concave in $\{b_{m,n}\}$, the feasible set defined by 
% the per-RU linear constraints is convex, and 
% Algorithm~\ref{alg:rmwf_p2} therefore converges to the global maximum 
% under this high-rate fronthaul regime.
% \end{remark}



\section{Numerical Results}\label{sec:numerical}

% {\color{red}{XXX}}

We consider an $M = 7$ cell hexagonal layout with an inter-site distance of $d = 300$~m, where each cell hosts a single RU equipped with $N$ antennas at a height of $30$~m.
% Assuming two users per cell, $K = 2M = 14$ users are uniformly distributed within their associated hexagonal cell at a height of $1.5$~m. 
In each channel realization, $K = 14$ users are dropped independently: each user is associated with a cell chosen uniformly at random and placed uniformly within that hexagonal cell at a height of $1.5$~m and at least $10$~m from its RU.
% An illustrative example of this layout is depicted in Fig.~\ref{fig:topology}.
Wireless channels are generated using the 3GPP 3D UMi standard channel model at a carrier frequency of $3.5$~GHz with a subcarrier spacing of $30$~kHz, including both pathloss and log-normal shadow fading. 
% We follow the convention of one resource element per channel realization, and each user transmits with a total power of $P = 10$~dBm distributed across a physical resource block of 168 resource elements (i.e., 12 subcarriers spanning 14 OFDM symbols), while the noise power spectral density is set to $-174$~dBm/Hz (i.e., noise power of $-129$~dBm for a $30$~kHz subcarrier spacing). 
% For the simulation setup, we consider a sum-power constraint over a physical resource block (PRB). 
% Specifically, a total power budget of $10$~dBm is uniformly normalized across the $168$ resource elements (i.e., $12$ subcarriers $\times$ $14$ OFDM symbols) of one PRB.
% This normalization yields an effective transmit power of $P_k = -12.25$~dBm per resource element. 
% The noise power spectral density is set to $-174$~dBm/Hz, corresponding to a per-RE noise power of $-129.2$~dBm for a $30$~kHz subcarrier spacing.
% Because the large-scale fading parameters remain constant across this PRB block, our evaluation is performed on a per-resource-element basis.
Each user is subject to a sum-power constraint of $10$~dBm over a physical resource block ($12$ subcarriers $\times$ $14$ OFDM symbols $=168$ resource elements (REs)), which is met with equal power per RE and thus gives $P_k = -12.25$~dBm. 
The per-RE noise power is $-129.2$~dBm, based on a $-174$~dBm/Hz density and a $30$~kHz subcarrier spacing.
Since the large-scale parameters are common across the block, we evaluate one RE per channel realization.
The modulation order $M_k$ for each user is drawn independently and uniformly at random from $\{4, 16, 64, 256\}$ in every channel realization. 
The antenna count $N$ and the fronthaul budget ($C_{\sf tot}$ or $C_m$) vary per figure and are stated in the corresponding caption.

% We consider an $M$ cell hexagonal layout with an inter-site distance of $d$~m, where each cell hosts a single RU equipped with $N$ antennas at a height of $30$~m. 
% Assuming two UEs per cell, $K=2M$ UEs are uniformly distributed within their associated hexagonal cell at a height of $1.5$~m. 
% An illustrative example of this layout is depicted in Fig.~\ref{fig:topology}.
% Wireless channels are generated using the 3GPP 3D urban-microcell (UMi) standard channel model at a carrier frequency of $3.5$~GHz with a subcarrier spacing of $30$~kHz, including both pathloss and log-normal shadow fading. 
% We follow the convention of one resource element per channel realization, and each UE transmits with a total power of $P$~dBm distributed across a physical resource block of 168 resource elements (i.e., 12 subcarriers spanning 14 OFDM symbols), while the noise power spectral density is set to $-174$~dBm/Hz (i.e., noise power of $-129$~dBm for a $30$~kHz subcarrier spacing). 
% The modulation order $M_k$ for each UE is drawn independently and uniformly at random from $\{4, 16, 64, 256\}$ in every channel realization. 

% Unless stated otherwise, we fix $M = 7$ RUs (hence $K = 2M = 14$ users), an
% inter-site distance $d = 300$~m, and a per-UE transmit power $P = 10$~dBm
% throughout this section; the antenna count $N$ and the fronthaul budget
% ($C_{\sf tot}$ for the sum constraint or $C_m$ for the per-RU constraint) are specified per figure.

We summarize the methods compared in our experiments as follows.
\begin{itemize}
    \item \textbf{Upper Bound}: The unquantized performance, obtained by setting the quantization error variance $q_{m,n} = 0$ for all $(m,n)$, equivalently $b_{m,n} \to \infty$. 
    This corresponds to a fronthaul with infinite capacity and provides a method-agnostic ceiling on the achievable sum GMI under the same channel and modulation realizations.
    
    % \item \textbf{RMWF}: The proposed Reverse Mercury/Waterfilling allocator with explicit finite-alphabet awareness. 
    % The mercury levels are evaluated using the MMSE lookup table, so the resulting allocation accounts for the saturation behavior of finite-alphabet inputs.

    \item \textbf{RMWF}:
    The proposed RMWF allocator with explicit finite-alphabet awareness. The MMSE lookup table (Fig.~\ref{fig:mmse_lut}) is constructed for $M_k \in \{4, 16, 64, 256\}$, with piecewise-linear interpolation between adjacent grid points at runtime. 
    % The proposed RMWF allocator with explicit finite-alphabet awareness. 
    % The required MMSE values $\mathrm{mmse}_{M_i}(\gamma)$ are evaluated through a precomputed lookup table, shown in Fig.~\ref{fig:mmse_lut}, that tabulates $\mathrm{mmse}_{M_i}(\gamma)$ over a dense SNR grid for each modulation order $M_i \in \{4, 16, 64, 256\}$. 
    % At runtime, MMSE values at arbitrary SNRs are obtained via piecewise-linear interpolation between adjacent grid points, so the resulting bit allocation accounts for the saturation behavior of finite-alphabet inputs.
    % The Panter-Dite quantizer constant $c = \pi\sqrt{3}/2$ is used.

    % \item \textbf{RMWF}: The proposed RMWF allocator with explicit finite-alphabet awareness. 
    % Both the MMSE function $\mathrm{mmse}_{M_i}(\gamma)$ and the mutual information $I_{M_i}(\gamma)$ are evaluated through precomputed lookup tables (LUTs) that tabulate each quantity over a dense SNR grid for every modulation order $M_i \in \{4, 16, 64, 256\}$. 
    % The MMSE LUT, shown in Fig.~\ref{fig:mmse_lut}, supplies the saturation scores in~\eqref{eq:score_update} that drive the bit update, while the GMI LUT supplies the per-iteration objective values used both for stopping (see Section~\ref{subsec:convergence}) and for reporting the final sum GMI. 
    % At runtime, both quantities at arbitrary SNRs are obtained via piecewise-linear interpolation between adjacent grid points, so the resulting bit allocation accounts for the saturation behavior of finite-alphabet inputs. 
    % The Panter--Dite quantizer constant $c = \pi\sqrt{3}/2$ is used.
    
    \item \textbf{RWF}:
    % A variant of the proposed allocator derived under the mismatched assumption of Gaussian inputs.
    % Despite the actual finite-alphabet signaling, this approach replaces the discrete MMSE function $\mathrm{mmse}_{M_i}(\gamma)$ with $1/(1 + \gamma)$.
    % Correspondingly, the constant is set to $c=1$, which is exact for Gaussian sources~\cite{CoverThomas}.
    The classical RWF approach derived under the sum-distortion minimization objective with the assumption of Gaussian sources~\cite{CoverThomas}. Unlike the proposed RMWF, this baseline allocates bits by minimizing the quantization error treating the inputs as Gaussian.
    

    \item \textbf{Equal-Bit}:
    % A modulation- and channel-agnostic baseline that allocates equal even integer budget only to the top half of the KLT coefficients corresponding to the largest eigenvalues. For instance, given $C_{\sf tot}=56$, $M=7$, and $N=4$, a budget of $b=2C_{\sf tot}/(MN)=4$ bits is equally assigned to each coefficient in the top half. It captures the performance of conventional scalar uniform quantization coupled with a rudimentary dimension reduction, without any adaptive bit allocation.
    % A modulation- and channel-agnostic baseline that assigns an equal even integer budget $b = 2C_{\sf tot}/(MN)$ only to the top half of the KLT coefficients corresponding to the largest eigenvalues, capturing conventional scalar uniform quantization with a rudimentary dimension reduction and no adaptive bit allocation.
    A modulation-agnostic baseline that retains only the $N/2$ dominant KLT coefficients per RU, assigning $2 \lfloor C_{\sf{tot}} / (MN) \rfloor$ bits to each retained coefficient and allocating any remaining bits in 2 bit increments to the strongest ones to exactly satisfy the budget.
\end{itemize}

% A scalar uniform quantizer is applied separately to the in-phase and quadrature components, so the real-valued RMWF/RWF allocations are rounded to non-negative even integers via the largest-remainder method; any residual budget mismatch is absorbed by $\pm 1$-bit adjustments to the largest entries, preserving the exact budget. 
% This rounding granularity also informs the stopping tolerance used in Section~\ref{subsec:convergence}.

% A scalar mid-rise uniform quantizer is applied at sample level, separately to the in-phase and quadrature components of each KLT coefficient: for $b_{m,n}/2$ bits per real dimension, the clipping interval $[-A, A]$ with $A = 3\sigma_{\rm r}$ and $\sigma_{\rm r}^2 = \lambda_{m,n}/2$ is partitioned into $2^{b_{m,n}/2}$ uniform bins and the sample is mapped to the corresponding bin-center reconstruction level, with the zero-bit case mapping to zero output. 
A scalar mid-rise uniform quantizer is applied at sample level, separately to the in-phase and quadrature components of each KLT coefficient, with subtractive dithering wrapped around it as in Section~\ref{subsec:cran_fa}: for $b_{m,n}/2$ bits per real dimension, the clipping interval $[-A, A]$ with $A = 3\sigma_{\sf r}$ and $\sigma_{\sf r}^2 = \lambda_{m,n}/2$ is partitioned into $2^{b_{m,n}/2}$ uniform bins and the sample is mapped to the corresponding bin-center reconstruction level, with the zero-bit case mapping to zero output. 
The $3\sigma_{\sf r}$ loading (i.e., $\kappa = 3$) yields the granular-regime constant $c = \kappa^2/3 = 3$ in $q_{m,n} = c\,\lambda_{m,n}\,2^{-b_{m,n}}$.
Since bits are assigned per real dimension, the real-valued RMWF/RWF allocations are rounded to non-negative even integers via the largest-remainder method, with any residual budget mismatch absorbed by $\pm 2$-bit adjustments to the largest entries to preserve the exact budget; this rounding granularity also informs the stopping tolerance used in Section~\ref{subsec:convergence}.



\subsection{GMI versus CCMI}\label{subsec:ccmi}

% \begin{figure}[t]     
% \centerline{\resizebox{0.8\columnwidth}{!}{\includegraphics{Figures/Topology/cran_topology_snapshot_matlab.eps}}}
%     \caption{
%     An illustrative example of the C-RAN topology for $M=7$ RUs and K=14 UEs, with an inter-site distance of $d=300$ m.
%     }\label{fig:topology}
% \end{figure}




Fig.~\ref{fig:ccmi} validates the use of the CCMI~\eqref{eq:mi_finite_def} 
as the RMWF design objective by directly comparing it with the actual 
achievable rate, the GMI~\eqref{eq:gmi}, for both (P1) and (P2).

After RMWF returns the bit allocation $\{b_{m,n}^\star\}$, the CU recomputes the quantization noise covariance $\mathbf{Q}^\star$ at this fixed allocation and redesigns the LMMSE receiver $\mathbf{W}_{\sf LMMSE}$ in~\eqref{eq:lmmse} accordingly. 
This two-stage use of $\mathbf{Q}$ is by design: the allocation stage requires the differentiable closed form $q_{m,n} = c\,\lambda_{m,n}\,2^{-b_{m,n}}$ to derive the fixed-point update, whereas the receiver stage, with $\{b_{m,n}^\star\}$ already fixed, can adopt the noise statistics, including the clipping contribution beyond the granular term. 
The CU evaluates these statistics from the signal statistics already available for detection, together with the loading factor $\kappa$, the bits it has assigned, and the locally generated dither, requiring no side information.

The GMI and CCMI curves are visually indistinguishable for both constraints across the entire range of $C_{\sf tot}$, so maximizing the CCMI~\eqref{eq:mi_finite_def} is empirically equivalent to maximizing the actual GMI~\eqref{eq:gmi} in our operating regime.
Guided by this validation, the remainder of this section adopts the 
following convention: RMWF performs bit allocation by maximizing 
the CCMI~\eqref{eq:mi_finite_def}, while all reported rates are 
evaluated as the actual GMI~\eqref{eq:gmi} obtained from Monte 
Carlo simulation of the true residual statistics.

\begin{figure}[t]     
\centerline{\resizebox{0.8\columnwidth}{!}{\includegraphics{Figures/CRAN/plot_cran_M7_K14_GMI13_vs_14_omni.eps}}}
    % \caption{
%     CCMI-based bit allocation under RMWF: the optimization objective (CCMI,
% markers) closely tracks the actual achievable rate (GMI, solid lines), validating
% the design choice. $M = 7$ RUs, $N = 16$ antennas per RU, and $K = 14$ users.}
\caption{CCMI-based bit allocation under RMWF: the optimization objective (CCMI, markers) closely tracks the achievable rate (GMI, solid lines), validating the design choice.
$N = 16$ antennas per RU, with $C_m = C_{\sf tot}/M$, $\forall m$, for (P2).}
\label{fig:ccmi}
\end{figure}


% \subsection{$N$-parallel Channel and Graphical Interpretation}\label{sec:toy}
% We begin with the toy example of $N$ independent parallel channels introduced at Section~\ref{sec:toy_fa}, where we additionally illustrate the mercury levels and vessel heights to highlight the necessity of finite-alphabet-aware bit allocation. 
% % In the following subsection, we then evaluate RMWF in the C-RAN environment described earlier to validate its effectiveness in a realistic deployment scenario.




% % Fig.~\ref{fig:toy_gmi_vs_b} compares the average sum GMI of the proposed RMWF against three baselines (RWF, Equal-bit, and the Upper Bound) on the toy $N$-parallel channel; details of the simulation setup are summarized in the caption.

% % % The proposed RMWF attains the highest average sum GMI across the entire bit-budget range. 
% % In general, the proposed RMWF demonstrates highly competitive average sum GMI performance.
% % Since most methods nearly close the gap to the Upper Bound starting from $b = 5$ bits per real scalar, the practically interesting regime is $b \in \{1, 2, 3, 4\}$, where the discrepancy between schemes is most pronounced. 
% % RMWF dominates throughout this range, with the largest gain occurring at $b = 2$: it achieves $33.66$ bits/symbol, compared to $28.4$ for RWF and $23.46$ for Equal-bit. 
% % In relative terms, RMWF improves over Equal-bit by $43.5\%$, while RWF improves over Equal-bit by only $21.03\%$. 
% % Equivalently, the finite-alphabet awareness of RMWF delivers an additional $18.54\%$ gain over its Gaussian counterpart RWF at this operating point.
% % These results confirm two key principles: (i) channel-SNR-aware bit allocation, as in the RMWF framework, substantially outperforms uniform allocation, and (ii) the finite-alphabet awareness of RMWF further amplifies this gain over its Gaussian counterpart.

% Fig.~\ref{fig:toy_G_vs_FA} compares the average sum GMI of the proposed RMWF against three baselines on the toy $N$-parallel channel.
% Note that we evaluate a special case consisting of two parallel channels to clearly highlight the importance of accounting for finite-alphabet inputs.
% Specifically, the first channel is configured with 4-QAM at a high transmit power-to-noise ratio of 20 dB, while the second channel employs 256-QAM at a relatively low 10 dB.
% As shown in Fig.~\ref{fig:toy_G_vs_FA}, the proposed RMWF achieves the highest average sum GMI among the considered schemes. 
% The largest performance gap occurs at $b=6$ bits per complex scalar, where RMWF provides a 7.01\% gain over Equal-bit and a 28.25\% gain over RWF.
% Conversely, the RWF, which relies on the Gaussian input assumption, exhibits severely degraded performance, falling even below that of the baseline equal-bit allocation.

% This performance degradation stems from the fundamental discrepancy between the Gaussian assumption and the actual characteristics of QAM signals, as discussed in Section~\ref{sec:toy_fa}.
% Allocation strategies based on the Gaussian assumption inherently tend to assign more quantization bits to channels with stronger signal power.
% However, in our scenario, the first channel (4-QAM, 20~dB) easily reaches saturation, as its MI approaches its theoretical maximum due to the sufficiently high SNR.
% In contrast, the second channel (256-QAM, 10~dB) operates far below saturation, meaning that additional quantization bits can significantly enhance its performance.

% \begin{figure}[t]     
% \centerline{\resizebox{0.9\columnwidth}{!}{\includegraphics{Figures/Toy/results_sigma2_1.0_fixedM4-256_P100-10_with_gaussian_N2.eps}}}
%      \caption{ Average sum GMI versus the number of bits per complex scalar $b$ for 2-parallel channels. 
%      The channel 1 with 4-QAM and $P_1/\sigma^2 = 20$dB.
%      The channel 2 with 256-QAM and $P_2/\sigma^2 = 10$dB.
%      }\label{fig:toy_G_vs_FA}
% \end{figure}

% \begin{figure}[t]     
% \centerline{\resizebox{1.0\columnwidth}{!}{\includegraphics{Figures/Toy/decomposition_snr_centric_matlab_b6.eps}}}
%     \caption{Graphical interpretation of RMWF and RWF, based on the high-SNR decomposition in~\eqref{eq:rmwf_sol}. 
%     The system environment is identical to that of Fig.~\ref{fig:toy_G_vs_FA}, and the given bit budget is $b = 6$ bits/complex scalar.
%     % For each channel, the allocated bits are obtained by filling water from the global water level over the mercury level, within the vessel determined by the post-quantization SNR. 
%     % The RWF mercury is depicted as zero, consistent with the high-SNR limit $\gamma\,\mathrm{mmse}_{\rm G}(\gamma) \to 1$.
%     }\label{fig:rmwf_rwf_interp}
% \end{figure}


% Fig.~\ref{fig:rmwf_rwf_interp} visualizes the decomposition
% in~\eqref{eq:rmwf_sol} for both schemes at $b=6$. 
% For each channel, the
% vessel height $\log_2\gamma_i^\star$ encodes the post-quantization SNR,
% the mercury level
% $\log_2(1/(\gamma_i^\star\mathrm{mmse}_{M_i}(\gamma_i^\star)))$ sits
% above the bottom, and the allocated bits are the clearance between the
% shared water level $\log_2 \mu^\star$ and the top of the mercury.

% The key contrast appears in channel~1 (4-QAM, 20\,dB). Under RMWF, the
% high SNR drives $\mathrm{mmse}_4(\gamma_1^\star)\to 0$, producing a
% thick mercury that blocks further bit deposition. RWF, treating the
% input as Gaussian, has $\gamma\,\mathrm{mmse}_{\rm G}(\gamma)\to 1$ and
% therefore zero mercury; it continues pouring bits into the
% already-saturated channel. The vessel heights themselves also differ,
% because $\gamma_i^\star$ depends on the final allocation: RMWF assigns
% fewer bits to channel~1 (smaller vessel) and more bits to channel~2
% (larger vessel), whereas RWF does the opposite. For channel~2 (256-QAM,
% 10\,dB), still far from saturation, the mercury is thin under both
% schemes; RMWF wins because it has redirected bits here where they
% materially improve detection. In short, RMWF discriminates by
% joint channel strength and modulation-dependent saturation,
% while RWF discriminates by signal power alone and wastes bits on
% resolved constellations.





\subsection{Average Sum GMI versus Bit Budget}

\begin{figure}[t]
    \centering
    % --- 첫 번째 그림 (위) ---
    \subfloat[\label{fig:8vs16_P1}]{
        \resizebox{0.8\columnwidth}{!}{\includegraphics{Figures/CRAN/TrueGMIresults_cran_M7_K14_P1_N8vsN16_omni.eps}}
    }
    \\ % <--- 핵심!! 이 줄바꿈이 위아래로 나누는 역할을 합니다.
    % --- 두 번째 그림 (아래) ---
    \subfloat[\label{fig:8vs16_P2}]{
        \resizebox{0.8\columnwidth}{!}{\includegraphics{Figures/CRAN/TrueGMIresults_cran_M7_K14_P2_N8vsN16_omni.eps}}
    }
    
%     \caption{
%         Average sum GMI versus total fronthaul capacity $C_{\sf tot}$ in an
% $M = 7$ cell hexagonal layout with an inter-site distance of $d = 300$~m, serving
% $K = 14$ users, each with a transmit power of $P = 10$~dBm, for $N \in \{8, 16\}$
% antennas per RU. (a) Sum fronthaul constraint. (b) Per-RU fronthaul constraint,
% where the total capacity is equally distributed among the RUs (i.e.,
% $C_m = C_{\sf tot}/M$, $\forall m$).
%     }
\caption{Average sum GMI versus total fronthaul capacity $C_{\sf tot}$ for
$N \in \{8, 16\}$ antennas per RU. (a) Sum fronthaul constraint. (b) Per-RU
fronthaul constraint with $C_m = C_{\sf tot}/M$, $\forall m$. }
    \label{fig:8vs16}
\end{figure}



% {\color{red}{XXX}}
Fig.~\ref{fig:8vs16} plots the average sum GMI as a function of $C_{\sf tot}$ for $N \in \{8, 16\}$ antennas per RU, with Fig.~\ref{fig:8vs16_P1} corresponding to the sum fronthaul constraint and Fig.~\ref{fig:8vs16_P2} to the per-RU constraint. 
Across both constraints and the entire range of $C_{\sf tot}$, the proposed RMWF outperforms both RWF and Equal-Bit. 
At the tightest budget $C_{\sf tot} = 112$ with $N = 16$, RMWF attains a maximum gain of approximately $484.84\%$ ($479.74\%$) over RWF (Equal-Bit) under the sum constraint and $422.84\%$ ($422.84\%$) under the per-RU constraint.

To trace the source of these gains, we first examine how the two baselines turn the bit budget into GMI.
Equal-Bit and RWF share the same design principle: both rank the KLT
coefficients by variance and remain blind to the constellation, differing only in that RWF activates coefficients adaptively through a water level, whereas Equal-Bit rigidly splits the budget over the top half of the eigenvalues. 
This yields a budget-dependent crossover in Fig.~\ref{fig:8vs16}: Equal-Bit outperforms RWF at tight budgets by concentrating bits on the strongest half, but is overtaken beyond $C_{\sf tot} = 560$ and $C_{\sf tot} = 1344$ for $N = 8$ and $N = 16$ respectively, once its excluded coefficients become the bottleneck while RWF keeps activating new ones.
Fixing the active set a priori is thus limiting; yet even RWF's adaptive selection still trails RMWF across the entire range.

RMWF overcomes both limitations simultaneously.
Because it allocates bits to directly maximize the finite-alphabet GMI, its selection is governed jointly by coefficient strength and the modulation-dependent saturation.
As a result, it neither fixes the active set in advance like Equal-Bit, nor, like RWF, wastes bits on coefficients whose marginal contribution to the GMI has already saturated, but instead flexibly selects, in every budget regime, the coefficients that contribute most to the GMI. 
This GMI-driven, finite-alphabet-aware selection is precisely what sustains
RMWF's advantage under tight fronthaul budgets.

% Two antenna counts $N \in \{8, 16\}$ are shown to indicate how the achievable
% ceiling scales with the array size---reflected by the gap between the two upper bounds---while the detailed dependence on $N$ at a fixed budget is examined in Section~\ref{subsec:gmi_vs_N}.





\subsection{Average Sum GMI versus the Number of Antennas}\label{subsec:gmi_vs_N}

\begin{figure}[t]
    \centering
    % --- 첫 번째 그림 (위) ---
    \subfloat[\label{fig:sweep_p1}]{
        \resizebox{0.8\columnwidth}{!}{\includegraphics{Figures/CRAN/TrueGMI_Sweep_M7_K14_Ctot896_P1_AvgSumGMI_omni.eps}}
    }
    \\ % <--- 핵심!! 이 줄바꿈이 위아래로 나누는 역할을 합니다.
    % --- 두 번째 그림 (아래) ---
    \subfloat[\label{fig:sweep_p2}]{
        \resizebox{0.8\columnwidth}{!}{\includegraphics{Figures/CRAN/TrueGMI_Sweep_M7_K14_Ctot896_P2_AvgSumGMI_omni.eps}}
    }
\caption{Average sum GMI versus the number of antennas per RU $N$ at a fixed total fronthaul capacity $C_{\sf tot} = 896$. 
(a) Sum fronthaul constraint. (b) Per-RU fronthaul constraint with $C_m = C_{\sf tot}/M = 128$.}
    \label{fig:sweep}
\end{figure}

Fig.~\ref{fig:sweep} plots the average sum GMI as a function of the number of antennas per RU $N$ at a fixed total fronthaul capacity $C_{\sf tot} = 896$, with Fig.~\ref{fig:sweep_p1} and Fig.~\ref{fig:sweep_p2} corresponding to the sum and per-RU fronthaul constraints, respectively. 
The upper bound, which is independent of any allocation strategy, increases monotonically with $N$, reflecting the additional spatial degrees of freedom available at the CU.

% The budget $C_{\sf tot} = 896$ is a deliberate reference point: as Fig.~\ref{fig:8vs16} shows, it is the budget at which both RMWF and RWF reach the upper bound at $N = 8$. Since increasing $N$ spreads the same budget over $MN = 7N$ coefficients, $N = 8$ separates two regimes.
% For $N \leq 8$ the budget is more than sufficient, so essentially every useful coefficient is quantized finely; for $N > 8$ the budget no longer suffices, and the way bits are allocated becomes decisive. Indeed, for $N \leq 8$ both RMWF and RWF rise steeply with $N$ and closely follow the upper bound, and they diverge only once the budget-limited regime $N > 8$ is entered.

% The budget $C_{\sf tot}=896$ is a deliberate reference point: as Fig.~\ref{fig:8vs16} shows, it is where both RMWF and RWF reach the upper bound at $N=8$. 
% In Fig.~\ref{fig:sweep}, since larger $N$ spreads the same budget over $MN=7N$ coefficients, $N=8$ marks the transition from a budget-rich regime ($N\le8$), where every useful coefficient is quantized finely and both schemes track the upper bound, to a budget-limited regime ($N>8$), where the allocation becomes decisive and the schemes diverge.


The budget $C_{\sf tot}=896$ is a deliberate reference point. 
As Fig.~\ref{fig:8vs16} shows, it provides exactly enough bits for both RMWF and RWF to reach the upper bound at $N=8$.
However, increasing $N$ spreads this fixed budget across more coefficients ($MN=7N$), making the allocation critical. Fig.~\ref{fig:sweep} fixes $C_{\sf tot}=896$ to explicitly capture this transition. 
It reveals a clear shift from a budget-rich regime ($N\le8$) where both schemes track the bound, to a budget-limited regime ($N>8$) where the schemes diverge.

The proposed RMWF settles into a near-flat plateau without any subsequent degradation once the budget-limited regime is entered, after gaining approximately $7.68$ and $7.31$~bits/symbol from $N = 4$ to $N = 12$ under the sum and per-RU constraints, respectively. 
The activation rule in \eqref{eq:bit_update_p1}--\eqref{eq:mu_closed} (analogously \eqref{eq:bit_update_p2}--\eqref{eq:mu_closed_p2} for the per-RU case) explains why: only coefficients with $c\lambda_{m,n} \cdot S_{m,n} > \mu$ enter the active set $\CMcal{A}$, so the extra coefficients introduced by larger $N$, carrying small eigenvalues, simply fall below the water level and are left inactive.
RMWF thus keeps concentrating the budget on the same dominant coefficients regardless of $N$, and the plateau is the imprint of this correct selection rather than an allocator-side limitation. 
The residual gap to the upper bound, which integrates all eigenmodes including the weak ones that RMWF excludes, widens only gradually with $N$.

In contrast, RWF degrades almost immediately past $N = 8$: ranking coefficients by variance alone under a Gaussian model, it cannot concentrate the scarce budget where it is needed, and its sum GMI falls monotonically thereafter.
Equal-Bit, on the other hand, keeps climbing until $N=14$ before degrading.
As Fig.~\ref{fig:8vs16} shows, Equal-Bit at $N=16$ is not fully saturated at $C_{\sf tot}=896$. 
Consequently, its peak in Fig.~\ref{fig:sweep} emerges at $N=14$, where $896$ bits suffice to hit a per-$N$ ceiling that stays below the upper bound due to top-half selection.
Beyond $N=14$, $b=2C_{\sf tot}/(MN)$ becomes so small that even the dominant coefficients are quantized too coarsely.
Lacking any mechanism to isolate the informative ones, Equal-Bit lets this distortion outweigh the spatial gain, losing about $24.52$~bits/symbol from its $N=14$ peak by $N=32$.

These behaviors carry a clear operational message: only RMWF turns additional antennas into a sustained benefit under a fixed fronthaul budget.
RMWF holds its GMI on a near-flat plateau as $N$ grows, whereas RWF and Equal-Bit convert antenna over-provisioning into a GMI loss. 
In short, RMWF can scale the antenna count while the competing schemes cannot.
This scalability is an increasingly consequential property as antenna counts grow faster than fronthaul capacity in modern C-RAN deployments.




\subsection{CDF of Per-User GMI}

\begin{figure}[t]     
\centerline{\resizebox{0.8\columnwidth}{!}{\includegraphics{Figures/CRAN/ecdf_cran_M7_N16_K14_Ctot224_AllMethods_LineAndMarkers_omni.eps}}}
\caption{CDF of per-user GMI for the sum (P1) and per-RU (P2) constraints;
$N = 16$ and $C_{\sf tot} = 224$.}
\label{fig:cdf}
\end{figure}

Fig.~\ref{fig:cdf} depicts the cumulative distribution function (CDF) of per-user GMI under the sum (P1) and per-RU (P2) fronthaul constraints. 
As shown in Fig.~\ref{fig:cdf}, under both Equal-Bit and RWF, approximately $50\%$ of the per-user GMI values fall below $0.35$~bits/symbol, and the vast majority remain below $4$~bits/symbol. 
This indicates that adaptive but Gaussian-based allocation provides little per-user benefit over Equal-Bit at this tight budget.

In contrast, the proposed RMWF achieves substantially higher per-user GMI. 
Under the sum constraint (P1), $50\%$ of the users obtain at least $3.71$~bits/symbol, and the $80$th percentile reaches $5.99$~bits/symbol. 
A similar trend is observed under the per-RU constraint (P2), where $50\%$ of the users obtain at least $3.25$~bits/symbol and the $80$th percentile reaches $5.68$~bits/symbol.

These results indicate that, regardless of whether a sum or per-RU fronthaul constraint is imposed, directly maximizing the finite-alphabet GMI is far more effective than minimizing reconstruction distortion under a Gaussian assumption.
This advantage manifests not only in the average sum GMI but also in the per-user GMI distribution, confirming that GMI-driven, finite-alphabet-aware bit allocation is critical from both system-level and user-level perspectives.
Furthermore, the RMWF advantage is somewhat larger under the sum constraint, since RMWF redistributes bits across RUs to lift more users.
The Gaussian-based RWF fails to exploit this flexibility, as its per-user distribution stays essentially unchanged between the two constraints.




\subsection{Convergence}
\label{subsec:convergence}

% \begin{figure}[t]     
% \centerline{\resizebox{0.9\columnwidth}{!}{\includegraphics{Figures/CRAN/Convergence_M7_N32_K14_breal1_GMI.eps}}}
%      \caption{
%      Average sum GMI versus iteration index $t$ under $M = 7$, $N = 32$, $K = 14$, $b = 2$, averaged over $S = 5000$ channel snapshots. 
%      Both Algorithm~\ref{alg:rmwf_p1} (P1, sum fronthaul) and Algorithm~\ref{alg:rmwf_p2} (P2, per-RU fronthaul) rise monotonically from the equal-bit initialization and reach a plateau within $T = 5$--$8$ iterations.
%      }\label{fig:GMIconv}
% \end{figure}

% \begin{figure}[t]     
% \centerline{\resizebox{0.9\columnwidth}{!}{\includegraphics{Figures/CRAN/Convergence_M7_N32_K14_breal1_Diff.eps}}}
%      \caption{
%      Per-element average bit-iterate change $\Delta^{(t)}$ defined in~\eqref{eq:bit_change} versus iteration index $t$, under the same configuration as Fig.~\ref{fig:GMIconv}.
%      The stopping tolerance $\epsilon = 0.3$ used in implementation is indicated; the small saturated-regime residual visible after $t \approx 10$ does not affect the rounded bit allocation.
%      }\label{fig:bitDifferConv}
% \end{figure}

% \begin{figure}[t]     
% \centerline{\resizebox{0.9\columnwidth}{!}{\includegraphics{Figures/CRAN/TrueConvergence_M7_N32_K14_breal1_GMI.eps}}}
%      \caption{
     % Average sum GMI versus iteration index $t$ under $M = 7$, $N = 32$, $K = 14$, $b = 2$, averaged over $S = 5000$ channel snapshots. 
     % Both Algorithm~\ref{alg:rmwf_p1} (P1, sum fronthaul) and Algorithm~\ref{alg:rmwf_p2} (P2, per-RU fronthaul) rise monotonically from the equal-bit initialization and reach a plateau within $T = 5$--$8$ iterations.
%      }\label{fig:GMIconv}
% \end{figure}

% \begin{figure}[t]     
% \centerline{\resizebox{0.9\columnwidth}{!}{\includegraphics{Figures/CRAN/TrueConvergence_M7_N32_K14_breal1_Diff.eps}}}
%      \caption{
%      Per-element average bit-iterate change $\Delta^{(t)}$ defined in~\eqref{eq:bit_change} versus iteration index $t$, under the same configuration as Fig.~\ref{fig:GMIconv}.
%      The stopping tolerance $\epsilon = 0.3$ used in implementation is indicated; the small saturated-regime residual visible after $t \approx 10$ does not affect the rounded bit allocation.
%      }\label{fig:bitDifferConv}
% \end{figure}

\begin{figure}[t]     
\centerline{\resizebox{0.8\columnwidth}{!}{\includegraphics{Figures/CRAN/TrueConvergence_M7_N32_K14_breal1_Combined_omni.eps}}}
%      \caption{
%      Average sum GMI versus iteration index $t$ under $M = 7$, $N = 32$,
% $K = 14$, and $C_{\sf tot} = 448$, averaged over $S = 5000$ channel snapshots. Both
% Algorithm~\ref{alg:rmwf_p1} (P1, sum fronthaul) and Algorithm~\ref{alg:rmwf_p2}
% (P2, per-RU fronthaul) rise monotonically from the equal-bit initialization and
% reach a plateau within $T = 5$--$8$ iterations. The lower panel shows the
% per-element average bit-iterate change $\Delta^{(t)}$ defined in
% \eqref{eq:bit_change} versus iteration $t$, under the same configuration; the stopping
% tolerance $\epsilon = 0.3$ is indicated, and the small saturated-regime residual
% visible after $t \approx 10$ does not affect the rounded bit allocation.}
\caption{Average sum GMI versus iteration index $t$ under $N = 32$ and $C_{\sf tot} = 448$, averaged over $S = 5000$ channel snapshots. 
Both Algorithm~\ref{alg:rmwf_p1} (P1) and Algorithm~\ref{alg:rmwf_p2} (P2) rise monotonically from the uniform initialization and reach a plateau within $T = 5$--$8$ iterations. 
The lower panel shows the per-element average bit-iterate change $\Delta^{(t)}$ defined in \eqref{eq:bit_change}; the stopping tolerance $\epsilon = 0.3$ is indicated, and the small saturated-regime residual visible after $t \approx 10$ does not affect the rounded bit allocation.}
\label{fig:GMIconv}
\end{figure}

A formal convergence proof is not pursued here; we instead report empirical evidence.
Fig.~\ref{fig:GMIconv} reports the convergence behavior over iterations. Its upper panel shows that the average sum GMI rises monotonically from the uniform initialization and plateaus within $T=5$--$8$ iterations for both (P1) and (P2).
The lower panel shows the per-element average bit change
\begin{align}
\Delta^{(t)} \triangleq \frac{1}{S}\sum_{s=1}^{S}
\frac{\|\mathbf{b}_s^{(t)} - \mathbf{b}_s^{(t-1)}\|_1}{MN},
\label{eq:bit_change}
\end{align}
averaged over $S=5000$ snapshots; 
the residual decays rapidly.
% during the first 10 iterations, with small post-plateau fluctuations due to near-zero-marginal-gain reshuffling that does not affect the objective. 
In implementation, we use the per-element-average stopping rule
$\Delta^{(t)}<\epsilon$ rather than $\|\cdot\|_\infty$ in Algorithms~\ref{alg:rmwf_p1}--\ref{alg:rmwf_p2}: norms on $\mathbb{R}^{MN}$ are equivalent, and the per-element form admits a size-independent reading, namely stabilization to within $\epsilon$ bits per coefficient on average. 
Because the real-valued allocation is subsequently rounded to even integers, sub-bit refinements have no effect on the quantized output, motivating $\epsilon=0.3$, below the rounding step. As the lower panel of Fig.~\ref{fig:GMIconv} shows, this threshold is met within $T\approx 6$ iterations for both (P1) and (P2), consistent with the GMI plateau.

% A formal convergence proof is not pursued here; we instead report empirical
% evidence. Fig.~\ref{fig:GMIconv} shows that the average sum GMI rises
% monotonically from the equal-bit initialization and plateaus within
% $T=5$--$8$ iterations for both (P1) and (P2). Fig.~\ref{fig:bitDifferConv}
% shows the per-element average bit change
% \begin{align}
% \Delta^{(t)} \triangleq \frac{1}{S}\sum_{s=1}^{S}
% \frac{\|\mathbf{b}_s^{(t)} - \mathbf{b}_s^{(t-1)}\|_1}{MN},
% \label{eq:bit_change}
% \end{align}
% averaged over $S=5000$ snapshots; the residual decays rapidly during the
% first ten iterations, with small post-plateau fluctuations due to
% near-zero-marginal-gain reshuffling that does not affect the objective.

% In implementation we use the per-element-average stopping rule
% $\Delta^{(t)}<\epsilon$ rather than $\|\cdot\|_\infty$ in
% Algorithms~\ref{alg:rmwf_p1}--\ref{alg:rmwf_p2}: norms on
% $\mathbb{R}^{MN}$ are equivalent, and the per-element form admits a
% size-independent reading (``stabilized to within $\epsilon$ bits per
% coefficient on average''). Because the real-valued allocation is
% subsequently rounded to even integers, sub-bit refinements have no
% effect on the quantized output, motivating $\epsilon=0.3$, below the
% rounding step. As Fig.~\ref{fig:bitDifferConv} shows, this threshold is
% met within $T\approx 8$ iterations for both (P1) and (P2), consistent
% with the GMI plateau.




\section{Conclusion}
\label{sec:conclusion}

% This paper has addressed the fronthaul bit allocation
% problem in uplink C-RAN under finite-alphabet inputs, with
% the end-to-end mutual information at the DU as the design
% objective. By applying the I-MMSE relation on the
% bit allocation side of the communication chain, we derived
% a stationary condition whose fixed-point iteration admits
% a closed-form update at every step. The converged
% allocation decomposes into a per-coefficient vessel height
% determined by the post-quantization SNR, a shared water
% level enforcing the fronthaul budget, and a finite-alphabet
% mercury level that quantifies the saturation state of each
% user's decoder. The resulting scheme, termed RMWF, generalizes classical reverse
% waterfilling through an explicit saturation correction,
% reduces exactly to it in the Gaussian limit, and completes
% a four-way correspondence between Gaussian and
% finite-alphabet inputs across the power and bit allocation
% sides of the communication chain. A single algorithmic
% structure handles both the sum and per-RU fronthaul
% provisioning regimes; the two formulations differ only in
% whether a single global water level or $M$ RU-local water
% levels govern the redistribution. 
% Numerical results under a
% 3GPP 3D UMi C-RAN deployment with mixed modulation orders
% confirmed that the proposed allocator substantially
% outperforms uniform and Gaussian-based allocations under
% tight fronthaul budgets, remains robust under antenna
% scaling where competing schemes degrade, and approaches the
% unquantized upper bound as the budget grows. The proposed
% iteration empirically converges within three to five
% rounds across all tested configurations. 


% We addressed fronthaul bit allocation for uplink C-RAN under finite-alphabet inputs, taking the end-to-end MI at the DU as the system objective.
We addressed fronthaul bit allocation for uplink C-RAN under finite-alphabet inputs, taking the post-LMMSE finite-alphabet GMI at the CU as the system objective.
Applied on the bit allocation side, the I-MMSE relation yields a fixed-point iteration termed RMWF.
Its converged allocation decomposes into a per-coefficient vessel height, a shared water level, and a finite-alphabet mercury level that captures decoder saturation.
Notably, the same structure handles both sum and per-RU fronthaul regimes, using either a single global water level or $M$ RU-local levels to govern the redistribution.
% The I-MMSE relation, applied on the bit allocation side, yields a fixed-point iteration, termed RMWF, whose converged allocation decomposes into a per-coefficient vessel height, a shared water level, and a finite-alphabet mercury level capturing decoder saturation; the same structure handles both sum and per-RU fronthaul regimes, with either a single global water level or $M$ RU-local levels governing the redistribution.
% RMWF generalizes classical reverse waterfilling through this saturation correction, reduces exactly to it in the Gaussian limit, and completes a four-way correspondence between Gaussian and finite-alphabet inputs across the power and bit allocation sides.
% RMWF generalizes classical reverse waterfilling through this saturation correction, recovers the latter under Gaussian inputs at high SNR, and completes a four-way correspondence between Gaussian and finite-alphabet inputs across the power and bit allocation sides.
RMWF generalizes classical RWF through this saturation correction and completes a four-way correspondence between Gaussian and finite-alphabet inputs across the power and bit allocation sides.
Numerical results on a 3GPP 3D UMi C-RAN with mixed modulation orders confirmed that RMWF substantially outperforms Equal-Bit and Gaussian baselines and remains robust under antenna scaling. 

% , and approaches the unquantized upper bound as the budget grows, with empirical convergence in 5 to 8 iterations.


% \appendices
% \section{Bounded-Support Argument for the High-Rate Inequality} \label{app:bounded_support}

% This appendix sketches the argument behind the inequality $I_{M_k}(\gamma_k) \le I_k^{\sf GMI}$ invoked in Remark~\ref{rem:gaussian_residual}, in the high-SNR/high-rate regime targeted by RMWF.

% Under the dithered construction, the post-LMMSE residual decomposes as $\eta_k = \eta_k^{\sf G} + \eta_k^{\sf B}$, where
% \begin{align}
%   \eta_k^{\sf G} \triangleq \mathbf{w}_k^{\sf H} \mathbf{n}, 
%   \qquad 
%   \eta_k^{\sf B} \triangleq \sum_{j \ne k} \mathbf{w}_k^{\sf H} \mathbf{h}_j x_j + \mathbf{w}_k^{\sf H} \mathbf{d}.
% \end{align}
% Then $\eta_k^{\sf G}$ is circularly Gaussian with variance $\sigma_{G,k}^2 \triangleq \sigma^2 \|\mathbf{w}_k\|^2$, and $\eta_k^{\sf B}$ has bounded support: each $|x_j|$ is bounded by the constellation peak, and $|e_{m,n}| \le \Delta_{m,n}/2$ by the Schuchman condition under~\eqref{eq:dithered_quant}. Hence there exists $r_k$, vanishing as the bit allocation grows, such that $|\eta_k^{\sf B}| \le r_k$ almost surely. Under the dithered construction, $\eta_k^{\sf G}$, $\eta_k^{\sf B}$, and $x_k$ are mutually independent.

% Let $d_k \triangleq |\rho_k|\sqrt{P_k}\, d_{\min}(\CMcal{X}_{M_k})$ denote the minimum spacing of the scaled constellation at the LMMSE output, and let $\gamma_k^{\sf G} \triangleq |\rho_k|^2 P_k / \sigma_{G,k}^2$ denote the Gaussian-only SNR. Since variances add under independence,
% \begin{align} \label{eq:appA_snr_gap}
%   \gamma_k^{\sf G} 
%   = \frac{|\rho_k|^2 P_k}{\sigma_{G,k}^2}
%   \;\ge\; 
%   \frac{|\rho_k|^2 P_k}{\sigma_{G,k}^2 + \mathbb{E}|\eta_k^{\sf B}|^2}
%   = \gamma_k.
% \end{align}

% When the high-rate condition $2 r_k < d_k$ holds, the bounded perturbation $\eta_k^{\sf B}$ cannot displace $\rho_k x_k$ across any decision boundary of the Gaussian nearest-neighbor decoder, leaving a guard margin of at least $d_k/2 - r_k$ to the nearest competing symbol. The pairwise error event is therefore contained in the Gaussian-only event $\{|\eta_k^{\sf G}| \ge d_k/2 - r_k\}$, equivalent to a Gaussian channel with effective squared minimum distance $(d_k - 2 r_k)^2$ in place of $d_k^2$. By the achievability of the GMI under i.i.d.\ codebooks~\cite{merhav1994information, ganti2000mismatched, martinez2009bit}, this yields
% \begin{align} \label{eq:appA_bound}
%   I_k^{\sf GMI} 
%   \;\ge\; 
%   I_{M_k}\!\left( \left(1 - \tfrac{2 r_k}{d_k}\right)^{\!2}\!\gamma_k^{\sf G} \right).
% \end{align}
% The right-hand side of~\eqref{eq:appA_bound} approaches $I_{M_k}(\gamma_k^{\sf G})$ as $r_k \to 0$ (i.e., as $\{b_{m,n}\}$ grow), and $I_{M_k}(\gamma_k^{\sf G}) \ge I_{M_k}(\gamma_k)$ by the monotonicity of $I_{M_k}(\cdot)$ in its argument~\cite{GuoShamaiVerdu} combined with~\eqref{eq:appA_snr_gap}. Hence
% \begin{align}
%   I_{M_k}(\gamma_k) \;\le\; I_k^{\sf GMI}
% \end{align}
% holds with arbitrarily small slack in the high-rate regime, justifying the use of $I_{M_k}(\gamma_k)$ as the design objective. For our setting ($M_k \in \{4, 16, 64, 256\}$ with $3\sigma$ loading), the condition $2 r_k < d_k$ is met at per-coefficient rates $b_{m,n} \gtrsim 5$--$9$ depending on the constellation, which coincides with the operating regime targeted by RMWF.



\bibliographystyle{IEEEtran}
\bibliography{ref_nfc}

\end{document}
